% Incognita paper P003. Rules: AGENTS.md, section 5; docs/writing.md (template preamble; at most
% 100 pages). Working notes go in notebook/, never in this file; what serves only the referees and
% the owner is in DOSSIER.md.
% One source, two builds (AGENTS.md 5.2), after `python3 scripts/certify.py bib`:
%   reader:    latexmk -pdf -interaction=nonstopmode -halt-on-error -cd papers/P003-transfer-setting/paper.tex
%   blueprint: latexmk -pdf -interaction=nonstopmode -halt-on-error -jobname=paper-blueprint
%                -usepretex='\def\IncognitaBlueprint{}' -cd papers/P003-transfer-setting/paper.tex
\documentclass[11pt]{article}
\usepackage[T1]{fontenc}
\usepackage[utf8]{inputenc}
\usepackage{lmodern}
\usepackage[a4paper,margin=2.5cm]{geometry}
\usepackage{amsmath,amssymb,mathtools}
\usepackage{booktabs}
\usepackage{algorithm}
\usepackage{algpseudocode}
\usepackage{xcolor}
\usepackage[colorlinks=true,linkcolor=blue!50!black,citecolor=blue!50!black,urlcolor=blue!50!black]{hyperref}
\usepackage[round]{natbib}
\usepackage{../../templates/paper/blueprint-macros}

% Macros of this paper: \newcommand only. Each symbol has one meaning in the whole paper.
% Navigation (docs/writing.md 2a): an anchor \dt{s:...} where a quantity is defined, and one macro
% per symbol built with \sy, used after its definition, so that the symbol links to it.
\newcommand{\dt}[1]{\hypertarget{#1}{}}
\newcommand{\sy}[2]{\hyperlink{#1}{#2}}
\newcommand{\E}{\mathbb{E}}
\newcommand{\PP}{\sy{s:PP}{\mathbb{P}}}
\newcommand{\KL}{\sy{s:KL}{\mathrm{KL}}}
\newcommand{\kl}{\sy{s:kl}{\mathrm{kl}}}
\newcommand{\N}{\mathbb{N}}
\newcommand{\ind}{\mathbf{1}}
\newcommand{\Pc}{\sy{s:Pc}{\mathcal{P}}}
\newcommand{\PcL}[1][L]{\sy{s:PcL}{\mathcal{P}_{#1}}}
\newcommand{\Sd}{\sy{s:Sd}{\mathcal{S}}}
\newcommand{\Gc}{\sy{s:Gc}{\mathcal{G}}}
\newcommand{\Bx}{\sy{s:Bx}{B}}
\newcommand{\ClL}[1][L]{\sy{s:Cl}{\mathcal{C}_{#1}}}
\newcommand{\Hs}{\sy{s:Hs}{\mathcal{H}}}
\newcommand{\Ber}{\sy{s:Ber}{\mathrm{Ber}}}
\newcommand{\Drift}{\sy{s:Drift}{\mathrm{Drift}}}
\newcommand{\one}{\sy{s:one}{\mathbf{1}_{K}}}
\newcommand{\Nla}{\sy{s:Nla}{N}^{\mathrm{la}}}
\newcommand{\pitr}{\sy{s:pitr}{\pi}^{\mathrm{tr}}}
\newcommand{\Np}{\sy{s:Np}{N'}}
\newcommand{\rhos}{\sy{s:Np}{\rho}}
\newcommand{\Dd}{\sy{s:D}{D}}
\newcommand{\PD}{\sy{s:PD}{P}^{D}}
\newcommand{\Q}{\sy{s:Q}{Q}}
\newcommand{\kap}{\sy{s:kappa}{\kappa}}
\newcommand{\Nc}{\sy{s:Nk}{N}}
\newcommand{\mean}{\sy{s:mu}{\mu}}
\newcommand{\gap}{\sy{s:mu}{\Delta}}
\newcommand{\Om}{\sy{s:Om}{\Omega'}}
\newcommand{\Up}{\sy{s:Om}{U'}}
\newcommand{\Xp}{\sy{s:Om}{X'}}
\newcommand{\mhat}{\sy{s:mhat}{\hat m'}}
\newcommand{\pis}{\sy{s:pis}{\pi}^{\sigma}}
\newcommand{\Pp}{\sy{s:Om}{P'}}
\newcommand{\Ncyc}{\sy{s:Ncyc}{N}^{\mathrm{cyc}}}
% Lean names of the results of Mathlib and LML used by name: printed (as a footnote) in the
% blueprint build only, never in the reader build (docs/writing.md, section 1).
\makeatletter
\newcommand{\leanfn}[1]{\ifincognita@blueprint\footnote{Lean: \nolinkurl{#1}.}\fi}
\makeatother

\title{The setting of transfer bandits with adaptively collected source data: size rules, the
consistency class $\mathcal{C}_L$, and the divergence decomposition}
\author{claude-opus-5-5, Incognita agents\thanks{Written by AI agents of the Incognita project
  (\url{https://github.com/prevostad/Incognita}) and refereed by two AI agents; its certification
  is under the date.}}
\date{2026-11-16}

\begin{document}
\maketitle

\begin{abstract}
Before a $K$-armed bandit game with rewards in $[0,1]$, a learner receives, for each arm $k$, the
first $N'_k$ values of an i.i.d.\ sequence drawn from a source law $\nu'_k$, where the sizes $N'_k$
are chosen by an unknown and possibly adaptive strategy, and the source law may differ from the
target law. We give a formal setting for this problem. A source strategy is a \emph{size rule}:
the event $\{N'=n\}$ depends only on the first $n_k$ values of each sequence and on independent
randomisation, that is, $N'$ is a stopping point of the $K$-parameter filtration of the source
sequences. The main result is the divergence decomposition: the Kullback--Leibler divergence of the
observations under two instances equals $\sum_k\E[N'_k]\,\mathrm{KL}(\nu'_k,\tilde\nu'_k)
+\sum_k\E[N_k(T)]\,\mathrm{KL}(\nu_k,\tilde\nu_k)$, for every size rule with $\E N'_k<\infty$, every
learner and every horizon, and it fails for a size map that looks one value ahead. The identity is
known when the source is a bandit game, and the proof combines classical facts: ignorability,
Wald's identity and the chain rule. With two arms, the bounded size rules are exactly the counts of
stopped bandit games; with three arms or more, a classical cyclic example shows that some size
rules are not, even in law. We also give witnesses and exclusions for the drift conditions and
for the consistency class.
\end{abstract}

\section{Introduction}\label{sec:intro}

In transfer learning for stochastic bandits, a learner facing $K$ arms with unknown reward laws
$\nu_1,\dots,\nu_K$ on $[0,1]$ also has data from a related \emph{source}: samples of laws
$\nu'_1,\dots,\nu'_K$ whose means are close to the target means, $\mean_k-\mean'_k<L_k$ for known
bounds $L_k$. When these samples were collected by another agent, their numbers $\Np_k$ are not
fixed in advance: the other agent chose them, possibly adaptively, by looking at the values it had
already seen. \citet{PMRM26} study this problem and prove a problem-dependent lower bound for a
class $\ClL$ of consistent learners, quantified over ``every prior strategy''; their lower bound
is stated for sources collected by a bandit game.

A setting for this problem must say which source strategies are allowed, what the learner sees, and
how a change of measure on the pair (source, online game) behaves: the lower bound of
\citet{PMRM26}, like the classical ones, needs the divergence between the laws of the observations
under two instances to depend on the source only through the expected sizes. This paper gives the
setting and proves that property for a broad class of strategies. Our main result
(Theorem~\ref{thm:main}) is the identity
\begin{equation}\label{eq:intro-decomposition}
\KL(\PP,\tilde\PP)=\sum_{k=1}^K\E[\Np_k]\,\KL(\nu'_k,\tilde\nu'_k)
+\sum_{k=1}^K\E[\Nc_k(T)]\,\KL(\nu_k,\tilde\nu_k),
\end{equation}
where $\PP$ and $\tilde\PP$ are the laws of the source data and of the first $T$ rounds under two
transfer instances $(\nu,\nu')$ and $(\tilde\nu,\tilde\nu')$, $\Nc_k(T)$ is the number of pulls of arm
$k$, and the expectations are under the first instance. It holds for every \emph{size rule} with
$\E \Np_k<\infty$, every learner and every horizon $T$, with no constant and no asymptotics; it
fails for a size map that looks one value ahead. Deterministic sizes and bandit games on the
source, stopped or not, are size rules (Lemma~\ref{lem:prior-game},
Proposition~\ref{prop:stopped-games}).

The identity~\eqref{eq:intro-decomposition} is known for bandit-game sources observed through their
ordered history \citep[Lemma~22]{PMRM26}, and follows in general from ignorability and Wald's
identity (Section~\ref{sec:related}). What this paper adds is the class of size rules,
the case of unbounded sizes with $\E \Np_k<\infty$, the source data as the learner's observation,
and an explicit size map for which the identity fails. The class is strictly larger than that of
stopped bandit games as soon as $K\ge3$, by a classical cyclic example (Section~\ref{sec:related}).

\paragraph{Our contributions} We summarize our main contributions below:
\begin{itemize}
\item A formal setting (Section~\ref{sec:setting}): size rules (Definition~\ref{def:size-rule}),
  source data (Definition~\ref{def:source-data}), transfer learners (Definition~\ref{def:learner}),
  the law of the game (Definition~\ref{def:game}), the drift conditions
  (Definition~\ref{def:drift}) and the class $\ClL$ (Definition~\ref{def:consistency}), with a
  discussion of the choices the source paper leaves open (Section~\ref{sec:reading}).
\item The divergence decomposition~\eqref{eq:intro-decomposition} for every size rule with
  $\E \Np_k<\infty$, and an explicit size map for which it fails (Theorem~\ref{thm:main}); the
  fundamental inequality of the change of measure follows (Corollary~\ref{cor:fundamental}). As
  said above, the identity itself is known or classical; we contribute its statement for all size
  rules, including unbounded ones, and the exclusion.
\item Witnesses and exclusions for each definition (Section~\ref{sec:properties}): bandit games,
  stopped or not, are size rules and the look-ahead map is not (Lemmas~\ref{lem:prior-game}
  and~\ref{lem:lookahead}, Proposition~\ref{prop:stopped-games}(a)); with two arms every bounded
  size rule is a stopped bandit game, and with three arms or more a classical cyclic size rule is
  not, even in law (Proposition~\ref{prop:stopped-games}(b), (c)); the relations between the drift
  conditions (Proposition~\ref{prop:drift}); $\ClL$ is antitone in $L$
  (Proposition~\ref{prop:classes}), contains every source-ignoring consistent algorithm
  (Proposition~\ref{prop:witness}), and excludes the learner that follows the source
  (Proposition~\ref{prop:exclusion}).
\end{itemize}

\paragraph{Outline} Section~\ref{sec:related} discusses related work. Section~\ref{sec:setting}
formalizes the setting. Section~\ref{sec:main} states the main result and
Section~\ref{sec:properties} the basic properties of the definitions. Section~\ref{sec:overview}
explains how the proofs work, Section~\ref{sec:examples} gives examples, and
Section~\ref{sec:open} concludes with open questions. The proofs are in the appendices.

\section{Related work and context}\label{sec:related}

\paragraph{Transfer and offline data in bandits} \citet{SJ12} study bandits with historical
observations collected before the game, and \citet[Section~2]{CL24} online bandits with a fixed
number of possibly biased offline samples per arm. In the transfer model of
\citet[Section~3]{PMM25}, the numbers of source samples are deterministic, from laws at a known
bounded discrepancy from the target laws; \citet[Section~3.2]{PMRM26} let them be chosen by an
unknown, possibly adaptive strategy, independent of the target, and define the class $\ClL$ of
Definition~\ref{def:consistency}. For best-arm identification, \citet[Section~2]{AJSS23} consider
offline samples collected by an unknown adaptive policy on the online instance.

\paragraph{Change of measure in bandits} The divergence decomposition of a bandit history,
$\KL=\sum_k\E[\Nc_k(T)]\KL(\nu_k,\tilde\nu_k)$, is \citet[Lemma~1]{KCG16} as an inequality on
events and \citet[(8)]{GMS19} as an identity, with data processing \citep[Lemma~1]{GMS19}. For a
sampling rule stopped at a stopping time, the expected log-likelihood ratio follows by Wald's
identity \citep{W44}, as in \citet[Appendix~A, (18)]{KCG16}.

\paragraph{Ignorability and missing data} A likelihood does not depend on a data-dependent choice
of what is observed when that choice uses only the observed values: this is the ignorability of
data missing at random \citep{R76}, made precise as ``everywhere missing at random'' by
\citet[Definition~2 and Theorem~1]{SGJC13}, and stated by \citet[Theorem~1]{FDS18}, for finitely
many variables, with the observed data a stopping-set $\sigma$-algebra. The condition
$\{\Np=n\}\in\Gc_n$ of Definition~\ref{def:size-rule} implies missingness at random with boxes as
observed sets, and the look-ahead map of Theorem~\ref{thm:main}(b) is not missing at random. In
sequential analysis this is the stopping-rule principle; see also \citet{BW88}.

\paragraph{Multiparameter processes and sequential observation} A size rule is a stopping point of
the $K$-parameter filtration $(\Gc_n)_{n\in\N^K}$; see also \citet{CW75}. A bandit strategy is an
optional increasing path over such a filtration \citep{M86,KM98}, and \citet[Section~1.3]{MV81}
call a stopping point \emph{predictable} when a policy reaches it, as a stopped bandit game does.
On $\N^2$ every stopping point is predictable under the conditional independence condition
$(\mathrm{F}4)$ \citep[Theorem~4]{MV81}; for three independent Markov chains a cyclic stopping point
is not \citep[Section~5.2]{MV81}; see also \citet{KS81}, \citet[Proposition~3]{M83} and \citet{W81}.
The cyclic pattern goes back to Berry's stable function that is not sequential
(\citealp[Proposition~II.2.2]{Ber76}; \citealp{Cur03}). In missing data, observing one variable at a time, chosen
from the values already seen, is the randomised monotone missingness of \citet{RG97}, who state
that some ignorable mechanisms are not of this form; see also \citet{GR97}.
Proposition~\ref{prop:stopped-games} gives the versions of these facts needed here.

\paragraph{Closest related work} \citet[Lemma~22]{PMRM26} prove~\eqref{eq:intro-decomposition} for
a bandit-game source whose ordered history the learner sees; it does not give
Theorem~\ref{thm:main}, since size rules need not be bandit games and our learner sees less.
\citet[(10)]{AJSS23} write the analogous expansion for offline data with the online laws; it does
not give Theorem~\ref{thm:main}, where the laws differ and the strategies need not be sequential.
Section~\ref{sec:main} compares them in detail, with Table~\ref{tab:comparison}.

\paragraph{Gap with the literature} For bounded size rules, Theorem~\ref{thm:main}(a) follows in
about a page from classical facts: ignorability gives Lemma~\ref{lem:source-law}, Wald's identity
the source term (Proposition~\ref{prop:source-kl}), and the chain rule the online term. We claim no
new idea for the identity; we write this route for size rules that may be unbounded, with
$\E \Np_k<\infty$, give a size map for which it fails, and give the setting itself.

\section{Problem formulation}\label{sec:setting}

\paragraph{Notational convention} Logarithms are natural; $\N:=\{0,1,2,\dots\}$,
$[K]:=\{1,\dots,K\}$; in $[0,\infty]$, $0\cdot\infty:=0$. Rounds are numbered from $1$. Inequalities
between vectors of $\N^K$ or $(0,\infty)^K$ are componentwise.

\subsection{Laws and instances}

\begin{definition}[{Laws on $[0,1]$}]\label{def:laws}
\dt{s:Pc}$\mathcal P$ is the set of Borel probability measures on $[0,1]$, and $m(P):=\int x\,P(dx)$
the mean of $P\in\Pc$; $\lambda$ is the uniform law on $[0,1]$,
\dt{s:Ber}$\mathrm{Ber}(p)$ the Bernoulli law of mean $p$, and $\delta_x$ the Dirac mass at
$x\in[0,1]$. For probability measures $P,Q$ on a measurable space,
\dt{s:KL}$\mathrm{KL}(P,Q):=\int\log\frac{dP}{dQ}\,dP$ if $P\ll Q$, and $\mathrm{KL}(P,Q):=\infty$
otherwise; \dt{s:kl}$\mathrm{kl}(p,q):=\KL(\Ber(p),\Ber(q))$.
\end{definition}

The integral defining $\KL(P,Q)$ exists in $[0,\infty]$: the negative part of the log-likelihood
ratio is integrable (Lemma~\ref{lem:kl-neg} in Appendix~\ref{sec:proofs}). So $\KL$ equals the
divergence of the Lean library Mathlib (Section~\ref{sec:overview})\leanfn{InformationTheory.klDiv},
which is $\infty$ when $P\not\ll Q$ or when the log-likelihood ratio is not integrable.

\begin{definition}[Instances]\label{def:instance}
\uses{def:laws}
Fix $K\ge2$. A target instance is $\nu=(\nu_1,\dots,\nu_K)\in\Pc^K$, with means\dt{s:mu} $\mean_k:=m(\nu_k)$,
$\mean^\star:=\max_k\mean_k$ and gaps $\gap_k:=\mean^\star-\mean_k$; arm $k$ is suboptimal if $\mean_k<\mean^\star$.
A source instance is $\nu'\in\Pc^K$, with means $\mean'_k:=m(\nu'_k)$. A transfer instance is a pair
$(\nu,\nu')\in\Pc^K\times\Pc^K$.
\end{definition}

\subsection{The source phase}\label{sec:setting-source}

The source is one probability space: for each arm an i.i.d.\ sequence (a \emph{stack}) of values
from its source law, and an independent randomisation for the strategy.

\begin{definition}[Source space]\label{def:source-space}
\uses{def:laws, def:instance}
\dt{s:Om}$\Om:=[0,1]^{\N}\times[0,1]^{[K]\times\N_{\ge1}}$ with the product $\sigma$-algebra and the
coordinates $\Up=(\Up_i)_{i\in\N}$ (randomisation) and $\Xp_{k,j}$ (the $j$-th value of the stack of
arm $k$). For $\nu'\in\Pc^K$, $\Pp_{\nu'}:=\lambda^{\otimes\N}\otimes\bigotimes_{k\in[K],j\ge1}\nu'_k$,
and $\E_{\nu'}$ is the expectation under $\Pp_{\nu'}$. For $n\in\N^K$, the box of $n$ is
\dt{s:Bx}$B(n):=\{(k,j):k\in[K],\,1\le j\le n_k\}$, and
\dt{s:Gc}$\mathcal G_n:=\sigma(\Up,\,\Xp_{k,j}:(k,j)\in\Bx(n))$.
\end{definition}

So under $\Pp_{\nu'}$ the variables $\Up$ and $\Xp_{k,j}$ are independent, $\Xp_{k,j}$ has law $\nu'_k$,
and $\Gc_n\subset\Gc_m$ when $n\le m$.

\begin{definition}[Size maps, size rules, source strategies]\label{def:size-rule}
\uses{def:source-space}
A size map is a measurable map\dt{s:Np} $\Np=(\Np_1,\dots,\Np_K):\Om\to\N^K$. It is a size rule if
\begin{equation*}
\{\Np=n\}\in\Gc_n\qquad\text{for every }n\in\N^K.
\end{equation*}
A source strategy is a sequence $\rhos=(\Np^{(T)})_{T\ge1}$ of size rules, $\Np^{(T)}$ being used at
horizon $T$.
\end{definition}

A size rule is a fixed map, which does not know $\nu'$; the law of $\Np$ under $\Pp_{\nu'}$ depends
on $\nu'$. The decision to stop at the sizes $n$ uses only the values in the box of $n$ and the
randomisation: $\Np$ is a stopping point of the $K$-parameter filtration $(\Gc_n)_{n\in\N^K}$.

\begin{definition}[Source data]\label{def:source-data}
\uses{def:source-space, def:size-rule}
\dt{s:Sd}$\mathcal S:=\bigsqcup_{n\in\N^K}\Sd_n$ with $\Sd_n:=[0,1]^{\Bx(n)}$ (a countable disjoint
union, with the sum $\sigma$-algebra); an element is $s=(n,x)$ with $x=(x_{k,j})_{(k,j)\in\Bx(n)}$,
and $n_k(s):=n_k$. For a size map $\Np$, the source data are\dt{s:D}
\begin{equation*}
\Dd:=\big(\Np,(\Xp_{k,j})_{(k,j)\in\Bx(\Np)}\big):\Om\to\Sd,
\end{equation*}
and\dt{s:PD} $\PD_{\nu',\Np}$ is the law of $\Dd$ under $\Pp_{\nu'}$. For $s$ with $n_k(s)\ge1$,\dt{s:mhat}
$\mhat_k(s):=n_k(s)^{-1}\sum_{j\le n_k(s)}x_{k,j}$ is the source empirical mean of arm $k$.
\end{definition}

$\Dd$ is measurable: on the event $\{\Np=n\}$ it equals $(n,(\Xp_{k,j})_{(k,j)\in\Bx(n)})$. The space
$\Sd$ is standard Borel.

\subsection{Learners and the law of the game}\label{sec:setting-game}

\begin{definition}[Bandit algorithms and transfer learners]\label{def:learner}
\uses{def:source-data}
Let \dt{s:Hs}$\mathcal H_t:=([K]\times[0,1])^{t}$ (with $\Hs_0$ a point). A bandit algorithm is a
sequence $\sigma=(\sigma_t)_{t\ge1}$ of Markov kernels from $\Hs_{t-1}$ to $[K]$. A transfer learner
is a sequence $\pi=(\pi_t)_{t\ge1}$ of Markov kernels from $\Sd\times\Hs_{t-1}$ to $[K]$; for
$s\in\Sd$, $\pi(s):=(\pi_t(s,\cdot))_{t\ge1}$ is a bandit algorithm. The source-ignoring learner of a
bandit algorithm $\sigma$ is\dt{s:pis} $\pis_t(s,h):=\sigma_t(h)$.
\end{definition}

A transfer learner is a bandit algorithm whose initial state is the source data. It knows $L$ and
$K$ (fixed before it is chosen), not the horizon.

\begin{definition}[Law of the game]\label{def:game}
\uses{def:instance, def:source-data, def:learner}
For a bandit algorithm $\sigma$, $\nu\in\Pc^K$ and $T\in\N$,\dt{s:Q} $\Q^T_{\sigma,\nu}$ is the law on $\Hs_T$
of $(A_1,Y_1,\dots,A_T,Y_T)$ when, for $t=1,\dots,T$, $A_t$ is drawn from $\sigma_t(H_{t-1})$ given
the history $H_{t-1}:=(A_i,Y_i)_{i<t}$ and $Y_t$ from $\nu_{A_t}$ given $(H_{t-1},A_t)$; that is,
$\Q^0_{\sigma,\nu}$ is the Dirac mass and
\begin{equation*}
\Q^t_{\sigma,\nu}:=\Q^{t-1}_{\sigma,\nu}\otimes(\sigma_t\otimes\nu),\qquad
(\sigma_t\otimes\nu)(h,\cdot):=\int\sigma_t(h,da)\,\delta_a\otimes\nu_a.
\end{equation*}
For $h\in\Hs_T$,\dt{s:Nk} $\Nc_k(T)(h):=\#\{t\le T:a_t=k\}$, and $\E_{\sigma,\nu}$ is the expectation under
$\Q^T_{\sigma,\nu}$. For a transfer learner $\pi$,\dt{s:kappa} $\kap^T_{\pi,\nu}(s):=\Q^T_{\pi(s),\nu}$ is a Markov
kernel from $\Sd$ to $\Hs_T$, and for a size map $\Np$ and $\nu'\in\Pc^K$ the law of the game at
horizon $T$ is\dt{s:PP}
\begin{equation*}
\mathbb P^T_{\nu',\Np,\nu,\pi}:=\PD_{\nu',\Np}\otimes\kap^T_{\pi,\nu}\quad\text{on }\Sd\times\Hs_T,
\end{equation*}
with coordinates $(\Dd,H_T)$; $\Np_k:=n_k(\Dd)$ and $\Nc_k(T):=\Nc_k(T)(H_T)$ on $\Sd\times\Hs_T$. For a
source strategy $\rhos$, $\E_{\nu',\rhos,\nu,\pi}$ is the expectation under
$\PP^T_{\nu',\Np^{(T)},\nu,\pi}$.
\end{definition}

The kernel $\kap^T_{\pi,\nu}$ is measurable in $s$ (finitely many composition products of kernels in
$(s,h)$). Given the source data, the online game depends neither on $\nu'$ nor on $\Np$: this is the
independence of the sizes and the target required by \citet[Section~3.2]{PMRM26}.

\subsection{Drift conditions and consistency}\label{sec:setting-class}

\begin{definition}[Drift conditions]\label{def:drift}
\uses{def:instance}
Let $L\in(0,\infty)^K$. The pair $(\nu,\nu')$ satisfies the drift condition
\dt{s:Drift}$\mathrm{Drift}_L$ if $\mean_k-\mean'_k<L_k$ for every $k$, and the two-sided strict
condition if $|\mean_k-\mean'_k|<L_k$ for every $k$. Following \citet[Section~4.2]{PMRM26},\dt{s:PcL}
\begin{equation*}
\mathcal P_L:=\big\{(\nu,\nu')\in\Pc^K\times\Pc^K:\ |\mean_k-\mean'_k|\le L_k\text{ for every }k\big\}.
\end{equation*}
We write \dt{s:one}$\mathbf 1_K:=(1,\dots,1)$.
\end{definition}

\begin{definition}[Consistency]\label{def:consistency}
\uses{def:instance, def:size-rule, def:learner, def:game, def:drift}
The class of \citet[Section~4.2]{PMRM26}, copied from the source: ``We define our class of
consistent algorithms as
\[
\mathcal{C}_{L}:=\left\{\pi\;\middle|\;\begin{array}[]{l}\forall(\nu,\nu^{\prime})\in\mathcal{P}([0,1])_{L},\\\forall\text{ prior strategy }\pi^{\prime},\\\forall\text{ suboptimal arm }k,\\\forall\alpha>0\end{array},\ \mathbb{E}_{\nu^{\prime},\pi^{\prime},\nu,\pi}[N_{k}(T)]=o(T^{\alpha})\right\}.\text{''}
\]
Formally, for $L\in(0,\infty)^K$: a transfer learner $\pi$ belongs to \dt{s:Cl}$\mathcal C_L$ if for
every $(\nu,\nu')\in\PcL$, every source strategy $\rhos$, every suboptimal arm $k$ of $\nu$ and every
$\alpha>0$,
\begin{equation*}
\lim_{T\to\infty}\frac{\E_{\nu',\rhos,\nu,\pi}[\Nc_k(T)]}{T^{\alpha}}=0.
\end{equation*}
A bandit algorithm $\sigma$ is consistent on $\Pc^K$ \citep[Section~4.1]{PMRM26} if for every
$\nu\in\Pc^K$, every suboptimal arm $k$ and every $\alpha>0$, $\E_{\sigma,\nu}[\Nc_k(T)]/T^{\alpha}\to0$
as $T\to\infty$.
\end{definition}

\section{The divergence decomposition}\label{sec:main}

For a fixed source, the law of the online history of a bandit algorithm satisfies the classical
decomposition $\KL(\Q^T_{\sigma,\nu},\Q^T_{\sigma,\tilde\nu})=\sum_k\E_{\sigma,\nu}[\Nc_k(T)]\KL(\nu_k,\tilde\nu_k)$
\citep[(8)]{GMS19}. The main result adds the source phase, for every size rule.

\begin{theorem}[Divergence decomposition]\label{thm:main}
\uses{def:laws, def:instance, def:source-space, def:size-rule, def:source-data, def:learner, def:game}
For every $K\ge2$, every size rule $\Np$, every transfer learner $\pi$, every $T\in\N$ and all pairs
$(\nu,\nu')$ and $(\tilde\nu,\tilde\nu')$ in $\Pc^K\times\Pc^K$ such that
\begin{equation*}
\E_{\nu'}[\Np_k]<\infty\qquad\text{for every arm }k,
\end{equation*}
write $\PP:=\PP^T_{\nu',\Np,\nu,\pi}$ and $\tilde\PP:=\PP^T_{\tilde\nu',\Np,\tilde\nu,\pi}$. Then:
\begin{enumerate}
\item[(a)] in $[0,\infty]$, with $0\cdot\infty:=0$,
\begin{equation}\label{eq:main}
\KL(\PP,\tilde\PP)=\sum_{k=1}^K\E_{\PP}[\Np_k]\,\KL(\nu'_k,\tilde\nu'_k)
+\sum_{k=1}^K\E_{\PP}[\Nc_k(T)]\,\KL(\nu_k,\tilde\nu_k),
\end{equation}
where $\E_{\PP}[\Np_k]=\E_{\nu'}[\Np_k]$ is the expected size of the source sample of arm $k$ under
$\nu'$ and $\Nc_k(T)$ is the number of pulls of arm $k$ in the first $T$ rounds;
\item[(b)] the hypothesis that $\Np$ is a size rule cannot be dropped. Let the look-ahead map
$N^{\mathrm{la}}:\Om\to\N^K$ be defined by\dt{s:Nla}
\begin{equation*}
N^{\mathrm{la}}_1:=1\ \text{if }\Xp_{1,2}>1/2,\qquad \Nla_1:=2\ \text{otherwise},\qquad
\Nla_k:=0\ \text{for }k\ge2.
\end{equation*}
It is a bounded size map and not a size rule. For all $p\ne q$ in $(0,1)$, every transfer learner
$\pi$, every $T\in\N$, every $\nu\in\Pc^K$ and all $\nu'_2,\dots,\nu'_K\in\Pc$, if $\nu'_1=\Ber(p)$,
$\tilde\nu'_1=\Ber(q)$, $\tilde\nu'_k=\nu'_k$ for $k\ge2$ and $\tilde\nu=\nu$, then the laws
$\PP:=\PP^T_{\nu',\Nla,\nu,\pi}$ and $\tilde\PP:=\PP^T_{\tilde\nu',\Nla,\nu,\pi}$ satisfy
\begin{equation*}
\KL(\PP,\tilde\PP)=2\,\kl(p,q),
\end{equation*}
while the right-hand side of~\eqref{eq:main} equals $(2-p)\,\kl(p,q)<2\,\kl(p,q)$.
\end{enumerate}
\end{theorem}
(see the proof in Appendix~\ref{sec:proofs}).

\paragraph{Constants and scope} There is no constant: (a) is an identity, for every horizon $T$ and
every pair of instances, with no asymptotics and no drift condition. The only hypothesis on the
strategy, besides being a size rule, is $\E_{\nu'}[\Np_k]<\infty$ under the first source instance;
nothing is assumed under the second one. We do not know whether (a) holds without this hypothesis:
the proof uses it to split the expected log-likelihood ratio into a positive and a negative part.
Every bandit game over $T'$ rounds has $\Np_k\le T'$, so the hypothesis holds for them. Part (b)
shows one size map outside the class; we do not claim that size rules are the largest class for
which (a) holds.

\paragraph{Interpretation} The source enters the divergence only through the expected sizes
$\E_{\nu'}[\Np_k]$, as if the sizes had been fixed at their means. Adaptivity of the source strategy
is invisible to a change of measure as long as the decision to stop reads only the values it keeps.
In (b) the size reveals the second value of the stack even when it is not kept, so the source data
carry the information of two samples for an expected size of $2-p$.

\paragraph{Comparison with the closest results} Table~\ref{tab:comparison} places
Theorem~\ref{thm:main} among the decompositions discussed in Section~\ref{sec:related}.
\begin{itemize}
\item \citet[Appendix~E, Lemma~22]{PMRM26} prove~\eqref{eq:main} when the source is a bandit game
  over $T'$ rounds: they write the prior and online games as one game of $T'+T$ rounds, whose
  ordered history the learner sees, and apply the divergence decomposition of bandit histories,
  under absolute continuity. It does not give Theorem~\ref{thm:main}: size rules need not be bandit
  games (Proposition~\ref{prop:stopped-games}(c)), and our learner sees the source data $\Dd$, a
  coarser observation than the ordered history.
\item \citet[Appendix~B.1, (10)]{AJSS23} write the expected log-likelihood ratio of offline and
  online data as $\sum_a(\E[N^{o}_a]+\E[\Nc_a])\KL(\nu_a,\tilde\nu_a)$, taking expectations ``with
  respect to all the randomness in the system''. Their offline data are sequential and have the
  online laws; it does not give Theorem~\ref{thm:main}, where the source laws differ from the target
  laws and the strategies need not be sequential.
\end{itemize}

\begin{table}[ht]
\centering
\caption{Divergence decompositions with a phase of data collected before the online game. ``Seq.''
means that the collection is a sequential sampling rule.}\label{tab:comparison}
\begin{tabular}{@{}p{0.30\textwidth}p{0.24\textwidth}p{0.18\textwidth}p{0.20\textwidth}@{}}
\toprule
Data before the game & Observation & Statement & Reference \\
\midrule
none & online history & identity & \citet[(8)]{GMS19} \\
one seq.\ rule at a stopping time, one law & ordered history & identity (Wald) & \citet[(18)]{KCG16} \\
offline, unknown adaptive policy, same laws online & offline and online data & asserted & \citet[(10)]{AJSS23} \\
bandit game over $T'$ rounds, source laws $\nu'$ & ordered prior history and online history & identity, $\nu\ll\tilde\nu$, $\nu'\ll\tilde\nu'$ & \citet[Lemma~22]{PMRM26} \\
size rule with $\E \Np_k<\infty$, source laws $\nu'$ & source data $\Dd$ and online history & identity in $[0,\infty]$; fails for a look-ahead map & Theorem~\ref{thm:main} \\
\bottomrule
\end{tabular}
\end{table}

\paragraph{The fundamental inequality} Combined with the data-processing inequality of
\citet[Lemma~1]{GMS19} (Lemma~\ref{cite:GMS19-Lemma1}), Theorem~\ref{thm:main}(a) gives the
inequality on which a lower bound rests, with the source term.

\begin{corollary}[Fundamental inequality with a source]\label{cor:fundamental}
\uses{def:laws, def:size-rule, def:source-data, def:learner, def:game, thm:main}
Under the hypotheses of Theorem~\ref{thm:main}(a), with $\PP$ and $\tilde\PP$ as there, for every
measurable $Z:\Sd\times\Hs_T\to[0,1]$,
\begin{equation*}
\sum_k\E_{\nu'}[\Np_k]\KL(\nu'_k,\tilde\nu'_k)+\sum_k\E_{\PP}[\Nc_k(T)]\KL(\nu_k,\tilde\nu_k)
\ge\kl\big(\E_{\PP}[Z],\E_{\tilde\PP}[Z]\big).
\end{equation*}
\end{corollary}
(see the proof in Appendix~\ref{sec:proof-cor}).

\section{Basic properties of the definitions}\label{sec:properties}

This section gives a witness and an exclusion for each definition of Section~\ref{sec:setting}. The
proofs are in Appendix~\ref{sec:proofs-properties}.

\paragraph{Size rules} Deterministic sizes and the counts of a prior bandit game are size rules;
the look-ahead map of Theorem~\ref{thm:main}(b) is not. A prior bandit game is a bandit game
played on the source stacks before the online game.

\begin{definition}[Prior bandit game]\label{def:prior-game}
\uses{def:source-space, def:size-rule}
Let $T'\in\N$ and let $\alpha_r:[0,1]^{\N}\times([K]\times[0,1])^{r-1}\to[K]$, $r=1,\dots,T'$, be
measurable.\dt{s:game0} On $\Om$ define $C_k(0):=0$ and, for $r=1,\dots,T'$,
\begin{equation*}
A'_r:=\alpha_r\big(\Up,(A'_i,Y'_i)_{i<r}\big),\qquad C_k(r):=C_k(r-1)+\ind\{A'_r=k\},\qquad
Y'_r:=\Xp_{A'_r,C_{A'_r}(r)}.
\end{equation*}
The counts of the prior game are $\Np_k:=C_k(T')$.
\end{definition}

This is the source of the lower bound of \citet[Theorem~2]{PMRM26}: a bandit game played over $T'$
rounds (possibly $T'=f(T)$) on the source stacks, by a strategy that uses its own observations. The
game may also stop early, at a time chosen with its own observations.

\begin{definition}[Stopped prior game]\label{def:stopped-game}
\uses{def:source-space, def:prior-game}
For a prior bandit game over $T'$ rounds,\dt{s:stopped} let $\mathcal F'_r:=\sigma\big(\Up,(A'_i,Y'_i)_{i\le r}\big)$
for $r=0,\dots,T'$. A stopping time of the game is a map $\tau:\Om\to\{0,\dots,T'\}$ with
$\{\tau=r\}\in\mathcal F'_r$ for every $r$, and the stopped counts are
$\Np_k:=C_k(\tau)=\sum_{r=0}^{T'}\ind\{\tau=r\}\,C_k(r)$.
\end{definition}

With $\tau\equiv T'$ the stopped counts are the counts of Definition~\ref{def:prior-game}.


\begin{lemma}[Witnesses of size rules]\label{lem:prior-game}
\uses{def:size-rule, def:prior-game}
(a) A constant map $\Np\equiv n^0\in\N^K$ is a size rule. (b) The counts $\Np$ of a prior bandit game
over $T'$ rounds (Definition~\ref{def:prior-game}) form a size rule, with $\Np_k\le T'$.
\end{lemma}

\begin{lemma}[Exclusion: looking ahead]\label{lem:lookahead}
\uses{def:size-rule}
The look-ahead map $\Nla:\Om\to\N^K$, $\Nla_1:=1$ if $\Xp_{1,2}>1/2$ and $\Nla_1:=2$ otherwise,
$\Nla_k:=0$ for $k\ge2$ (the map of Theorem~\ref{thm:main}(b)), is a bounded size map and not a
size rule.
\end{lemma}

\paragraph{Size rules and stopped bandit games} Is a size rule just a bandit game on the source
stacks, stopped when it chooses? With two arms, yes for bounded size rules; with three arms or more,
no, not even in law. Write\dt{s:ek} $e_k$ for the $k$-th unit vector of $\N^K$ and $|n|:=\sum_kn_k$.

\begin{proposition}[Size rules and stopped bandit games]\label{prop:stopped-games}
\uses{def:source-space, def:size-rule, def:prior-game, def:stopped-game}
\begin{enumerate}
\item[(a)] The stopped counts of a prior game over $T'$ rounds (Definition~\ref{def:stopped-game})
  form a size rule, with $|\Np|\le T'$.
\item[(b)] Let $K=2$ and let $\Np$ be a size rule with $T':=\sup_{\Om}|\Np|<\infty$. There are a
  prior game over $T'$ rounds and a stopping time $\tau$ of it whose stopped counts equal $\Np$ at
  every point of $\Om$.
\item[(c)] Let $K\ge3$ and $b_k:=\ind\{\Xp_{k,1}>1/2\}$ for $k=1,2,3$. The cyclic map
  $N^{\mathrm{cyc}}:\Om\to\N^K$,\dt{s:Ncyc}
  \begin{equation*}
  N^{\mathrm{cyc}}:=\begin{cases}
  e_1+e_2&\text{on }\{b_1=1,\,b_2=0\},\\
  e_2+e_3&\text{on }\{b_2=1,\,b_3=0\},\\
  e_3+e_1&\text{on }\{b_3=1,\,b_1=0\},\\
  e_1+e_2+e_3&\text{elsewhere},
  \end{cases}
  \end{equation*}
  is a size rule with $|\Ncyc|\le3$. There is no prior game, over any number $T'$ of rounds, with a
  stopping time whose stopped counts have, for every $\nu'\in\Pc^K$, the same law under $\Pp_{\nu'}$
  as $\Ncyc$.
\end{enumerate}
\end{proposition}

Since the sizes are a function of the source data, no stopped prior game has the same law of source
data as $\Ncyc$ for every $\nu'$ either: for $K\ge3$, size rules are strictly more general than
stopped bandit games, also for the learner, who sees only $\Dd$.

\paragraph{Sources of Proposition~\ref{prop:stopped-games}} Part (b) is a version of
\citet[Theorem~4]{MV81} that holds pointwise and for all source laws at once. Part (c) is a
classical example, and we claim nothing new in it beyond its writing for size rules. The pattern of
$\Ncyc$ is that of a function of \citet[Proposition~II.2.2]{Ber76}, the ``Gustave function'' in
the form of \citet{Cur03}, defined on three boolean arguments by $G(1,0,\bot)=G(\bot,1,0)=G(0,\bot,1)=1$, where $\bot$ is an argument not
read: the arguments read on these three inputs are the observed arms of $\Ncyc$ on its first three
cases. Berry shows that no argument can be evaluated first on every input, so that no sequential
program computes $G$; this is the argument of our proof, and the example of
\citet[Section~5.2]{MV81} for three Markov chains. With stacks of depth one ($\Np\le\one$), a size
rule is a mechanism of missing data that is missing at random (Lemma~\ref{lem:source-law}), and a
stopped prior game with counts at most $\one$ is a randomised monotone missingness scheme in the
sense of \citet{RG97}, who
state that some ignorable mechanisms are not of this form; see also \citet{GR97}. So part (c), in
law and with randomisation, is that phenomenon for size rules, a class between monotone schemes and
mechanisms missing at random.

\paragraph{Source data and learners} A transfer learner sees the source only through $\Dd$: it cannot
use a value that the size rule did not keep.

\begin{lemma}[The learner does not see unobserved values]\label{lem:learner-exclusion}
\uses{def:source-data, def:learner}
Let $\Np\equiv(1,0,\dots,0)$. There is no measurable $f:\Sd\to[0,1]$ with $f(\Dd)=\Xp_{1,2}$ on
$\Om$. Hence there is no transfer learner $\pi$ with $\pi_1(\Dd(\omega),\cdot)=\delta_{1}$ when
$\Xp_{1,2}(\omega)>1/2$ and $\pi_1(\Dd(\omega),\cdot)=\delta_2$ otherwise, for every $\omega\in\Om$.
\end{lemma}

\paragraph{Drift conditions} The three conditions of Definition~\ref{def:drift} differ, and the
next proposition says exactly when one implies another.

\begin{proposition}[Relations between the drift conditions]\label{prop:drift}
\uses{def:drift}
Let $L\in(0,\infty)^K$.
\begin{enumerate}
\item[(a)] The two-sided strict condition implies $\Drift_L$ and $(\nu,\nu')\in\PcL$.
\item[(b)] $\PcL\subset\{(\nu,\nu'):\Drift_{L'}\}$ for every $L'\in(0,\infty)^K$ with $L'_k>L_k$ for all $k$.
\item[(c)] $\{\Drift_L\}\subset\PcL$ if and only if $L_k\ge1$ for every $k$; and
  $\PcL\subset\{\Drift_L\}$ if and only if $L_k>1$ for every $k$. In particular, for
  $L\in(0,1)^K$ neither set contains the other.
\item[(d)] $\PcL=\Pc^K\times\Pc^K$ if and only if $L\ge\one$; $\Drift_L$ holds for every pair if and
  only if $L_k>1$ for every $k$.
\item[(e)] Witness and exclusion: every pair with $\nu'=\nu$ satisfies all three conditions for every
  $L$; the pair with $\nu_k=\delta_1$, $\nu'_k=\delta_0$ for every $k$ is in $\PcL[\one]$ and does not
  satisfy $\Drift_{\one}$.
\end{enumerate}
\end{proposition}

\paragraph{The consistency class} $\ClL$ shrinks when $L$ grows, stops changing at $L=\one$, and
lies between two explicit learners.

\begin{proposition}[Monotonicity of $\mathcal{C}_L$]\label{prop:classes}
\uses{def:consistency, def:drift}
(a) If $L\le L'$, then $\ClL[L']\subset\ClL$. (b) $\ClL=\ClL[\one]$ whenever $L\ge\one$, and
$\bigcap_{L\in(0,\infty)^K}\ClL=\ClL[\one]$: a transfer learner is in $\ClL$ for every $L$ if and
only if it satisfies the condition of Definition~\ref{def:consistency} with $\PcL$ replaced by
$\Pc^K\times\Pc^K$, that is, for every pair and every source strategy.
\end{proposition}

\begin{proposition}[Witness: source-ignoring learners]\label{prop:witness}
\uses{def:consistency, def:learner, def:game}
If a bandit algorithm $\sigma$ is consistent on $\Pc^K$, its source-ignoring learner $\pis$
belongs to $\ClL[\one]$, hence to $\ClL$ for every $L\in(0,\infty)^K$. The algorithm UCB
\citep{ACF02}, in the form formalised in the Lean library LML (Section~\ref{sec:overview}), with
parameter $c>1/2$ (pull arm $t$ at rounds
$t\le K$, then at round $t>K$ an arm maximising
$\hat\mu_a(t-1)+\sqrt{2c\log t/\Nc_a(t-1)}$, with $\hat\mu_a(t-1)$ and $\Nc_a(t-1)$ the empirical mean
and the count of arm $a$ after $t-1$ rounds and ties broken by a fixed rule) is consistent on
$\Pc^K$; in particular $\ClL\ne\emptyset$.
\end{proposition}

\begin{proposition}[Exclusion: following the source]\label{prop:exclusion}
\uses{def:consistency, def:learner, def:source-data}
Let \dt{s:pitr}$\pi^{\mathrm{tr}}$ be the transfer learner that, when $n_k(s)\ge1$ for every $k$,
pulls at every round the smallest arm $\hat k(s)$ maximising $\mhat_k(s)$, and otherwise pulls
arm $1$ at every round. Then $\pitr\notin\ClL$ for every $L\in(0,\infty)^K$, and the failure already
occurs on a pair with $\nu'=\nu$ and deterministic source sizes.
\end{proposition}

\subsection{Reading of the source}\label{sec:reading}

Six points of Definitions~\ref{def:learner}--\ref{def:consistency} are choices where
\citet{PMRM26} leave freedom.
\begin{itemize}
\item ``Prior strategy'' is not defined formally in \citet{PMRM26}: its Section~3.2 says that the
  counts may come from an online game, may depend on $T$ and on the prior observations, and are
  independent of the target. We read it as a source strategy (Definition~\ref{def:size-rule}): a
  class for which Theorem~\ref{thm:main}(a) holds, which contains the bandit games of
  \citet[Theorem~2]{PMRM26}, stopped or not (Lemma~\ref{lem:prior-game},
  Proposition~\ref{prop:stopped-games}(a)), and outside which (a) can fail
  (Theorem~\ref{thm:main}(b)). We do not claim that it is the largest such class. Restricting the
  strategies (to bandit games, or to deterministic sizes) can only enlarge $\ClL$ (Open
  question~\ref{oq:deterministic}).
\item \citet[Section~3.1]{PMRM26} assume a single optimal arm; the definition of $\ClL$ quantifies
  over all of $\mathcal P([0,1])_L$, and we keep it literally: every pair, ``suboptimal'' meaning
  $\mean_k<\mean^\star$.
\item The set $\PcL$ is closed and two-sided, while the drift condition of
  \citet[Section~3.2]{PMRM26} is strict, and the one-sided strict condition $\Drift_L$ is the one
  under which their upper bound holds (their Section~7 notes that the lower inequalities play no
  role there). This is why the three conditions are defined; Proposition~\ref{prop:drift} compares
  them.
\item Learners are anytime (Definition~\ref{def:learner}), while the source strategy may depend on
  $T$.
\item The online phase is a composition product (Definition~\ref{def:game}) rather than a game on
  stacks. Online rewards drawn from stacks $X_{k,j}\sim\nu_k$, the reward of round $t$ being
  $X_{A_t,\Nc_{A_t}(t)}$, give the same law of $(\Dd,H_T)$; this realisation is not used.
\item In the proof of the lower bound of \citet[Appendix~E]{PMRM26}, the learner of the online
  phase sees the whole ordered prior history (the actions, the values and the randomisation of the
  prior game). Our learner sees the source data $\Dd$ only: the sizes and the values of each arm, in
  their order within the arm. For a bandit-game source, $\Dd$ is a function of the ordered prior
  history, so every transfer learner is also a learner in their sense; a learner that uses the
  interleaving of the arms in the prior game is not a transfer learner here. The results of
  Sections~\ref{sec:main} and~\ref{sec:properties} do not depend on this choice except
  Theorem~\ref{thm:main}, whose observation is $(\Dd,H_T)$.
\end{itemize}

The propositions on $\ClL$ use only deterministic sizes and source-ignoring learners, so they hold
verbatim for the class defined with bandit-game sources only, or with the ordered prior history as
the learner's observation.

\section{How the proofs work}\label{sec:overview}

\paragraph{The decomposition} The law of the game is a composition product
$\PP=\PD_{\nu',\Np}\otimes\kap^T_{\pi,\nu}$ (Definition~\ref{def:game}). The chain rule for the
Kullback--Leibler divergence of composition products splits $\KL(\PP,\tilde\PP)$ into the divergence
of the source data and the average divergence of the online phase:
\begin{equation*}
\KL(\PP,\tilde\PP)=\underbrace{\KL\big(\PD_{\nu',\Np},\PD_{\tilde\nu',\Np}\big)}_{\text{source}}
+\underbrace{\int\KL\big(\kap^T_{\pi,\nu}(s),\kap^T_{\pi,\tilde\nu}(s)\big)\,\PD_{\nu',\Np}(ds)}_{\text{online}}.
\end{equation*}
The online term is the classical decomposition of bandit histories for each fixed $s$, integrated
in $s$ (Lemmas~\ref{lem:mixture}, \ref{lem:lml-history} and~\ref{lem:online-kl}). The source term
is handled in three steps:
\begin{enumerate}
\item \emph{Factorisation} (Lemma~\ref{lem:source-law}): on each slice $\{n\}\times\Sd_n$, the law
  of $\Dd$ is $g_n\cdot\bigotimes_{(k,j)\in\Bx(n)}\nu'_k$, with a function $g_n\in[0,1]$ that does not
  depend on $\nu'$. This is where the size-rule condition is used: $\{\Np=n\}$ is an event of
  $\Gc_n$, so integrating out the randomisation leaves a function of the kept values only.
\item \emph{Likelihood ratio}: since $g_n$ is the same under $\nu'$ and $\tilde\nu'$, the density of
  $\PD_{\nu',\Np}$ with respect to $\PD_{\tilde\nu',\Np}$ is the product of the per-sample ratios
  $\frac{d\nu'_k}{d\tilde\nu'_k}(x_{k,j})$.
\item \emph{Wald} (Lemma~\ref{lem:wald}): the expectation of a sum over the kept values of arm $k$
  is $\E_{\nu'}[\Np_k]$ times the expectation of one value, because the event $\{\Np_k\ge j\}$ is
  independent of $\Xp_{k,j}$.
\end{enumerate}
Proposition~\ref{prop:source-kl} combines them.

\paragraph{Where the difficulty lies} Each step is classical; the care is in the integrability.
The expected log-likelihood ratio is a sum of a positive and a negative part, and the negative
part of each per-sample ratio has integral at most $1/e$ (Lemma~\ref{lem:kl-neg}); the hypothesis
$\E_{\nu'}[\Np_k]<\infty$ makes the sum of these negative parts integrable, which justifies the split.
The case where $\nu'_k\not\ll\tilde\nu'_k$ for an arm that is sampled with positive probability is
treated separately (both sides are infinite). For (b), the chain rule holds for every size map; the
look-ahead map makes $\Dd$ a bijective image of the first two values of the stack, so the source
divergence is that of two samples.

\paragraph{Formal libraries} The setting is written to be formalised in the proof assistant Lean, on
top of Mathlib, its library of mathematics (\url{https://github.com/leanprover-community/mathlib4}),
and of LML, a Lean library of formalised results on bandits and sequential learning
(\url{https://github.com/LeanMachineLearning/LML}). A few textbook steps of the proofs use results of
these libraries by name: the chain rule for the Kullback--Leibler divergence, the divergence
decomposition of bandit histories and a regret bound of UCB.

\paragraph{Map of the appendices} Appendix~\ref{sec:cited} states the data-processing inequality
of \citet{GMS19}. Appendix~\ref{sec:proofs} proves Theorem~\ref{thm:main} and
Corollary~\ref{cor:fundamental}; Appendix~\ref{sec:proofs-properties} proves the results of
Section~\ref{sec:properties}; Appendix~\ref{sec:attribution} lists the sources of the techniques
and Appendix~\ref{sec:notation} the notation.

\section{Examples}\label{sec:examples}

This paper proposes no algorithm, so there is nothing to simulate; instead we compute the two sides
of~\eqref{eq:main} exactly on small instances.

\begin{example}[An adaptive rule computed by hand]\label{ex:hand}
Let $K\ge2$ and sample the source of arm $1$ once, and once more if the first value is $1$, with no
other arm sampled: $\Np_1:=1+\ind\{\Xp_{1,1}=1\}$ and $\Np_k:=0$ for $k\ge2$. This is a size rule:
$\{\Np=(1,0,\dots,0)\}=\{\Xp_{1,1}\ne1\}$ is in $\Gc_{(1,0,\dots,0)}$,
$\{\Np=(2,0,\dots,0)\}=\{\Xp_{1,1}=1\}$ is in $\Gc_{(2,0,\dots,0)}$, and the other events $\{\Np=n\}$
are empty. Let
$\nu'_1=\Ber(p)$, $\tilde\nu'_1=\Ber(q)$ and $\tilde\nu'_k=\nu'_k$ for $k\ge2$. The source data are
$((1,0,\dots,0),(0))$ with probability $1-p$, and $((2,0,\dots,0),(1,y))$ with probability
$p\cdot p^y(1-p)^{1-y}$ for $y\in\{0,1\}$. Hence
\begin{align*}
\KL\big(\PD_{\nu',\Np},\PD_{\tilde\nu',\Np}\big)
&=(1-p)\log\frac{1-p}{1-q}+p\log\frac pq+p\,\kl(p,q)\\
&=(1+p)\,\kl(p,q),
\end{align*}
and $\E_{\nu'}[\Np_1]=1+p$, as~\eqref{eq:main} says (with $T=0$).
\end{example}

\begin{example}[The look-ahead map]\label{ex:lookahead}
For the map $\Nla$ of Theorem~\ref{thm:main}(b), ``keep the second value only if it is $0$'', the
expected size is $\E_{\nu'}[\Nla_1]=2-p$ but the divergence is $2\,\kl(p,q)$: the size reveals
$\Xp_{1,2}$. With $p=0.3$ and $q=0.5$, $\kl(0.3,0.5)=0.08228$, so
\begin{equation*}
\KL\big(\PD_{\nu',\Nla},\PD_{\tilde\nu',\Nla}\big)=0.16457,\qquad
\E_{\nu'}[\Nla_1]\,\kl(p,q)=0.13988,
\end{equation*}
while the rule of Example~\ref{ex:hand} gives $(1+p)\kl(p,q)=0.10697$ on both sides.
\end{example}

\begin{example}[All the hypotheses at once]\label{ex:all}
Let $K=2$, $\nu=\nu'=(\Ber(2/3),\Ber(1/3))$, $\tilde\nu=\tilde\nu'=(\Ber(2/3),\Ber(3/4))$ and
$L=(0.1,0.1)$, so that $(\nu,\nu')$ satisfies $\Drift_L$, the two-sided strict condition and
$(\nu,\nu')\in\PcL$. Take the source strategy of round-robin prior games over
$T':=\lceil\log(T+1)\rceil$ rounds that pull the arms in the order $1,2,1,2,\dots$ (size rules with
$\Np_k\le T'$, Lemma~\ref{lem:prior-game}), and the source-ignoring UCB with $c=1$, which is in $\ClL$
(Proposition~\ref{prop:witness}). Then $\E[\Np_2]=\lfloor T'/2\rfloor$ and~\eqref{eq:main} reads
\begin{equation*}
\KL(\PP,\tilde\PP)=\big(\lfloor T'/2\rfloor+\E[\Nc_2(T)]\big)\,\kl\big(\tfrac13,\tfrac34\big).
\end{equation*}
\end{example}

\paragraph{Degenerate cases} For $\Np\equiv0$, $\Dd$ is a point and Theorem~\ref{thm:main}(a) is the
bandit decomposition of \citet[(8)]{GMS19}. For deterministic sizes $\Np\equiv n$, the source term is
$\sum_kn_k\KL(\nu'_k,\tilde\nu'_k)$. For $T=0$, $\Hs_0$ is a point and (a) is
Proposition~\ref{prop:source-kl}. If $\nu'_1=\Ber(1/2)$ and $\tilde\nu'_1=\delta_0$, then
$\KL(\nu'_1,\tilde\nu'_1)=\infty$, and both sides of (a) are infinite as soon as
$\Pp_{\nu'}(\Np_1\ge1)>0$; if arm $1$ is $\Pp_{\nu'}$-almost surely never sampled, its term is
$0\cdot\infty=0$. Ties in the means play no role in Theorem~\ref{thm:main}; in
Definition~\ref{def:consistency}, an instance with several optimal arms constrains only the arms
with $\mean_k<\mean^\star$.

\section{Open questions}\label{sec:open}

\paragraph{Conclusion}
We gave a formal setting for transfer bandits whose source sizes are chosen by an unknown adaptive
strategy, and showed that the change of measure behaves as with fixed sizes, the source entering
through $\E \Np_k$ only, for every size rule with $\E \Np_k<\infty$. The limitations: the lower and
upper bounds of \citet{PMRM26} are not proved here (Proposition~\ref{prop:drift} only compares the
sets of instances, not the bounds under each condition); the hypothesis $\E_{\nu'}[\Np_k]<\infty$ is
not shown to be necessary; and the choice of size rules as the class of source strategies, and of
$\Dd$ as the learner's observation, are modelling choices (Section~\ref{sec:reading}). For $K\ge3$
these choices matter: size rules are strictly more general than stopped bandit games
(Proposition~\ref{prop:stopped-games}(c)), and whether this changes the consistency class is part
of the question below.

\begin{openquestion}[Deterministic sizes against size rules]\label{oq:deterministic}
\extends{D1}
For $L\in(0,1)^K$, let $\mathcal C^{\rm det}_L$ be the class of Definition~\ref{def:consistency}
with the source strategies restricted to deterministic sizes ($\Np^{(T)}\equiv n(T)$ for sequences
$n(T)\in\N^K$), and $\mathcal C^{\rm game}_L$ the class with the strategies restricted to prior
bandit games (Definition~\ref{def:prior-game}) over $T'(T)$ rounds. Then
$\ClL\subset\mathcal C^{\rm game}_L\subset\mathcal C^{\rm det}_L$. Is either inclusion strict: is
there a transfer learner, consistent for every deterministic sequence of sizes, that fails against
an adaptive strategy, for instance one that keeps sampling the source of the optimal arm until its
empirical mean looks low? Theorem~\ref{thm:main} cannot separate the classes, since it sees the
strategy only through $\E \Np_k$; a separation needs a learner whose use of the source is valid for
fixed sizes only (confidence bounds that are not uniform in the sample size), and must keep
consistency on the closed set $\PcL$, including $\mean'_1+L_1=\mean_1$. The answer says whether the
class of a lower bound depends on how the source was collected.
\end{openquestion}

\appendix

\section{Useful lemmas}\label{sec:cited}

In this section we present the lemma from other papers that we use in our proofs: the
data-processing inequality, used in Corollary~\ref{cor:fundamental}.

\citet{GMS19} write $(\Gamma,\mathcal G)$ for the measurable space and $\mathbb P_1,\mathbb P_2$ for
the two laws; the domain of $Z$, printed $\Omega$ in the source, is $\Gamma$. We keep their
notation.

% Verbatim: literature/notes/GMS19.md, Lemma 1.
\begin{lemma}[{From \citet[Lemma~1]{GMS19}}]\label{cite:GMS19-Lemma1}
Consider a measurable space $(\Gamma,\mathcal G)$ equipped with two distributions $\mathbb P_1$ and
$\mathbb P_2$, and any $\mathcal G$--measurable random variable $Z:\Gamma\to[0,1]$. We denote
respectively by $\mathbb E_1$ and $\mathbb E_2$ the expectations under $\mathbb P_1$ and $\mathbb P_2$.
Then, $\KL(\mathbb P_1,\mathbb P_2)\ge\kl(\mathbb E_1[Z],\mathbb E_2[Z])$.
\end{lemma}

\section{Proof of Theorem~\ref{thm:main}}\label{sec:proofs}

To prove this theorem we use the chain rule of Section~\ref{sec:overview}; the following lemmas
control the two terms.

\paragraph{The source term} First a fact on divergences, then the three steps of
Section~\ref{sec:overview}.

\begin{lemma}[Negative part of the log-likelihood ratio]\label{lem:kl-neg}
\uses{def:laws}
If $P\ll Q$ with $r:=dP/dQ$, then
\begin{equation*}
\int(\log r)^-\,dP\le\frac1e,
\end{equation*}
so the integral defining $\KL(P,Q)$ exists in $[0,\infty]$.
\end{lemma}
(see the proof in Subsection~\ref{sec:proof-kl-neg}).

\begin{lemma}[Law of the source data]\label{lem:source-law}
\uses{def:source-space, def:size-rule, def:source-data}
Let $\Np$ be a size rule. For every $n\in\N^K$ there is a measurable $g_n:\Sd_n\to[0,1]$, which does
not depend on $\nu'$, such that for every $\nu'\in\Pc^K$ and every measurable $G\subset\Sd_n$,
\[
\PD_{\nu',\Np}(\{n\}\times G)=\int_Gg_n\,d\Big(\bigotimes_{(k,j)\in\Bx(n)}\nu'_k\Big).
\]
\end{lemma}
(see the proof in Subsection~\ref{sec:proof-source-law}). This is an instance of the
ignorability of a missing-at-random mechanism, which \citet[Theorem~1]{FDS18} state for finitely
many variables; here there are countably many, and an independent randomisation.

\begin{lemma}[Wald's identity for size rules]\label{lem:wald}
\uses{def:source-space, def:size-rule}
Let $\Np$ be a size rule, $\nu'\in\Pc^K$, $k\in[K]$ and $\varphi:[0,1]\to[0,\infty]$ measurable.
Then, in $[0,\infty]$,
\begin{equation*}
\E_{\nu'}\Big[\sum_{j=1}^{\Np_k}\varphi(\Xp_{k,j})\Big]=\E_{\nu'}[\Np_k]\int\varphi\,d\nu'_k.
\end{equation*}
If $\E_{\nu'}[\Np_k]<\infty$ and $\int|\varphi|\,d\nu'_k<\infty$ for a real $\varphi$, the same
identity holds with finite terms.
\end{lemma}
(see the proof in Subsection~\ref{sec:proof-wald}). The proof is Wald's, with a stopping point in
place of a stopping time.

\begin{proposition}[Divergence of the source data]\label{prop:source-kl}
\uses{def:laws, def:source-space, def:size-rule, def:source-data}
Let $\Np$ be a size rule and $\nu',\tilde\nu'\in\Pc^K$ with $\E_{\nu'}[\Np_k]<\infty$ for every $k$.
Then $\KL(\PD_{\nu',\Np},\PD_{\tilde\nu',\Np})=\sum_k\E_{\nu'}[\Np_k]\,\KL(\nu'_k,\tilde\nu'_k)$.
\end{proposition}
(see the proof in Subsection~\ref{sec:proof-source-kl}).

\paragraph{The online term} The law of the game is a mixture over the source data, and for each
value of the source data the online phase is a bandit game; Lemma~\ref{lem:lml-history} identifies
its law with that of LML, whose divergence decomposition we use.

\begin{lemma}[Mixture over the source]\label{lem:mixture}
\uses{def:game}
For a size map $\Np$, a transfer learner $\pi$, $\nu,\nu'\in\Pc^K$, $T\in\N$ and measurable
$f:\Sd\times\Hs_T\to[0,\infty]$,
$\E_{\PP^T_{\nu',\Np,\nu,\pi}}[f]=\int\E_{\pi(s),\nu}[f(s,\cdot)]\,\PD_{\nu',\Np}(ds)$. In particular
$\E_{\PP}[\Np_k]=\E_{\nu'}[\Np_k]$ and $\E_{\PP}[\Nc_k(T)]=\int\E_{\pi(s),\nu}[\Nc_k(T)]\,\PD_{\nu',\Np}(ds)$.
\end{lemma}
(see the proof in Subsection~\ref{sec:proof-mixture}).

\begin{lemma}[The online law as an LML history]\label{lem:lml-history}
\uses{def:laws, def:learner, def:game}
Let $\sigma$ be a bandit algorithm and $\nu\in\Pc^K$. Let $c:\mathbb R\to[0,1]$,
$c(y):=\min(\max(y,0),1)$, act on $([K]\times\mathbb R)^{t}$ coordinatewise on the rewards, let
$\bar\sigma_t:=\sigma_t\circ c$ (a Markov kernel from $([K]\times\mathbb R)^{t-1}$ to $[K]$), let
$\bar\nu_k$ be the image of $\nu_k$ by the inclusion $[0,1]\subset\mathbb R$, and let
$\iota_t:\Hs_t\to([K]\times\mathbb R)^{t}$ be the inclusion. Let $(A_t,Y_t)_{t\ge1}$ be a sequence of
random variables with values in $[K]\times\mathbb R$ such that, for every $t\ge1$, the conditional
law of $A_t$ given $H_{t-1}:=(A_i,Y_i)_{i<t}$ is $\bar\sigma_t(H_{t-1})$ and the conditional law of
$Y_t$ given $(H_{t-1},A_t)$ is $\bar\nu_{A_t}$. Then for every $T\in\N$, the law of $H_T$ is the
image of $\Q^T_{\sigma,\nu}$ by $\iota_T$.

In particular this holds for an algorithm--environment sequence of LML, read as follows. In LML the
rounds are numbered from $0$, and before each action the algorithm receives an observation; for
bandits the observations take values in a one-point space. A history of $t$ rounds of LML is a
sequence of $t$ triples (observation, action, reward), and forgetting the observations is a
measurable bijection $j_t$ from these histories onto $([K]\times\mathbb R)^{t}$, with a measurable
inverse. Take the algorithm of LML whose policy of index $t-1$ is $\bar\sigma_t\circ j_{t-1}$ (the
observation of the round being ignored) and the bandit environment of LML given by $\bar\nu$. Then
the actions and rewards of index $t-1$ of an algorithm--environment sequence of LML are $(A_t,Y_t)$
above, and the law of its history of $T$ rounds is the image of $\iota_T\#\Q^T_{\sigma,\nu}$ by
$j_T^{-1}$.
\end{lemma}
(see the proof in Subsection~\ref{sec:proof-lml-history}).

\begin{lemma}[Divergence of the online phase]\label{lem:online-kl}
\uses{def:laws, def:game}
For a size map $\Np$, a transfer learner $\pi$, $\nu,\tilde\nu,\nu'\in\Pc^K$ and $T\in\N$, with
$\PP:=\PP^T_{\nu',\Np,\nu,\pi}$,
\begin{equation*}
\int\KL\big(\kap^T_{\pi,\nu}(s),\kap^T_{\pi,\tilde\nu}(s)\big)\,\PD_{\nu',\Np}(ds)
=\sum_k\E_{\PP}[\Nc_k(T)]\KL(\nu_k,\tilde\nu_k).
\end{equation*}
\end{lemma}
(see the proof in Subsection~\ref{sec:proof-online-kl}). For a fixed source this is the classical
decomposition \citep[(8)]{GMS19}; no contribution here.

\begin{proof}[Proof of Theorem~\ref{thm:main}]
\proves{thm:main}
\uses{prop:source-kl, lem:online-kl, lem:mixture, lem:lookahead, def:laws, def:instance, def:source-space, def:size-rule, def:source-data, def:learner, def:game}
\emph{(a)} By the chain rule for the Kullback--Leibler divergence of composition products (a
textbook fact, formalised in LML\leanfn{InformationTheory.klDiv_compProd_eq_add_lintegral}; it
applies since $\Sd$ and $\Hs_T$ are standard Borel),
\begin{equation*}
\KL(\PP,\tilde\PP)=\KL(\PD_{\nu',\Np},\PD_{\tilde\nu',\Np})
+\int\KL(\kap^T_{\pi,\nu}(s),\kap^T_{\pi,\tilde\nu}(s))\,\PD_{\nu',\Np}(ds).
\end{equation*}
The first term is $\sum_k\E_{\nu'}[\Np_k]\KL(\nu'_k,\tilde\nu'_k)$ by
Proposition~\ref{prop:source-kl} (\emph{divergence of the source data}), and
$\E_{\nu'}[\Np_k]=\E_{\PP}[\Np_k]$ by Lemma~\ref{lem:mixture}; the second is
$\sum_k\E_{\PP}[\Nc_k(T)]\KL(\nu_k,\tilde\nu_k)$ by Lemma~\ref{lem:online-kl}. This is~\eqref{eq:main}.

\emph{(b)} The map $\Nla$ is a bounded size map and not a size rule by Lemma~\ref{lem:lookahead}.
The chain rule above holds for every size map (it does not use Proposition~\ref{prop:source-kl}).
Since $\tilde\nu=\nu$, $\kap^T_{\pi,\tilde\nu}=\kap^T_{\pi,\nu}$ and the second term vanishes, so
\begin{equation*}
\KL(\PP,\tilde\PP)=\KL\big(\PD_{\nu',\Nla},\PD_{\tilde\nu',\Nla}\big).
\end{equation*}
Under $\Pp_{\nu'}$ and under $\Pp_{\tilde\nu'}$, almost surely $\Xp_{1,1},\Xp_{1,2}\in\{0,1\}$,
$\Nla_1=1$ if and only if $\Xp_{1,2}=1$, and $\Dd=\Phi(\Xp_{1,1},\Xp_{1,2})$ with
\begin{equation*}
\Phi(a,1):=\big((1,0,\dots,0),(a)\big),\qquad \Phi(a,0):=\big((2,0,\dots,0),(a,0)\big).
\end{equation*}
The map $\Phi$ is injective on $\{0,1\}^2$, so $\PD_{\nu',\Nla}$ and $\PD_{\tilde\nu',\Nla}$ are the
images of $\Ber(p)^{\otimes2}$ and $\Ber(q)^{\otimes2}$ by $\Phi$, two laws on four atoms carried
bijectively; hence
\begin{equation*}
\KL\big(\PD_{\nu',\Nla},\PD_{\tilde\nu',\Nla}\big)=\KL\big(\Ber(p)^{\otimes2},\Ber(q)^{\otimes2}\big)=2\,\kl(p,q).
\end{equation*}
On the other side, $\E_{\nu'}[\Nla_1]=p+2(1-p)=2-p$, $\Nla_k=0$ for $k\ge2$, and
$\KL(\nu_k,\tilde\nu_k)=0$ for every $k$, so the right-hand side of~\eqref{eq:main} is
$(2-p)\,\kl(p,q)$. Since $p\ne q$ are in $(0,1)$, $0<\kl(p,q)<\infty$, and $p>0$ gives
$(2-p)\kl(p,q)<2\kl(p,q)$.
\end{proof}

\subsection{Proof of Corollary~\ref{cor:fundamental}}\label{sec:proof-cor}

\begin{proof}
\uses{thm:main, cite:GMS19-Lemma1, def:game, def:source-data}
Apply Lemma~\ref{cite:GMS19-Lemma1} on $(\Gamma,\mathcal G):=\Sd\times\Hs_T$ with its product
$\sigma$-algebra, $\mathbb P_1:=\PP$ and $\mathbb P_2:=\tilde\PP$ (probability measures by
Definition~\ref{def:game}): $\KL(\PP,\tilde\PP)\ge\kl(\E_{\PP}[Z],\E_{\tilde\PP}[Z])$. Replace
$\KL(\PP,\tilde\PP)$ by the right-hand side of Theorem~\ref{thm:main}(a), with
$\E_{\PP}[\Np_k]=\E_{\nu'}[\Np_k]$.
\end{proof}

\subsection{Proof of Lemma~\ref{lem:kl-neg}}\label{sec:proof-kl-neg}

\begin{proof}
\uses{def:laws}
The set $\{r=0\}$ has $P$-measure $0$, and $x(\log x)^-=x\log(1/x)\le1/e$ for $x\in(0,1)$, so
\begin{equation*}
\int(\log r)^-\,dP=\int r(\log r)^-\,dQ\le\sup_{0<x<1}x\log(1/x)=1/e.
\end{equation*}
So the integral of $\log r$ against $P$ exists in $(-\infty,\infty]$. It is nonnegative by Jensen's
inequality: $\int\log(1/r)\,dP\le\log\int r^{-1}\,dP=\log Q(r>0)\le0$.
\end{proof}

\subsection{Proof of Lemma~\ref{lem:source-law}}\label{sec:proof-source-law}

\begin{proof}
\uses{def:source-space, def:size-rule, def:source-data}
Let $Y_n:=(\Up,(\Xp_{k,j})_{(k,j)\in\Bx(n)})$, with values in $[0,1]^{\N}\times\Sd_n$. Since
$\Gc_n=\sigma(Y_n)$ and $\{\Np=n\}\in\Gc_n$, there is a measurable set $E_n$ with
$\{\Np=n\}=\{Y_n\in E_n\}$. On $\{\Np=n\}$, $\Dd=(n,(\Xp_{k,j})_{\Bx(n)})$, so
\begin{equation*}
\PD_{\nu',\Np}(\{n\}\times G)=\Pp_{\nu'}\big(Y_n\in E_n,\ (\Xp_{k,j})_{\Bx(n)}\in G\big).
\end{equation*}
Under $\Pp_{\nu'}$, $Y_n$ has law $\lambda^{\otimes\N}\otimes\bigotimes_{\Bx(n)}\nu'_k$ (independent
coordinates). By Tonelli, the probability is $\int_G g_n\,d(\bigotimes_{\Bx(n)}\nu'_k)$ with
\begin{equation*}
g_n(x):=\int\ind_{E_n}(u,x)\,\lambda^{\otimes\N}(du)\in[0,1],
\end{equation*}
measurable in $x$. The set $E_n$ is determined by the map $\Np$ ($Y_n$ is a surjective coordinate
projection), so $g_n$ depends only on $\Np$, not on $\nu'$.
\end{proof}

\subsection{Proof of Lemma~\ref{lem:wald}}\label{sec:proof-wald}

\begin{proof}
\uses{def:source-space, def:size-rule}
For $j\ge1$ let $\mathcal F^{k,j}:=\sigma(\Up,\,\Xp_{i,l}\ (i\ne k,\,l\ge1),\,\Xp_{k,l}\ (l<j))$. Then
\begin{equation*}
\{\Np_k\le j-1\}=\bigcup_{n:\,n_k\le j-1}\{\Np=n\}
\end{equation*}
is a countable union of events of $\Gc_n\subset\mathcal F^{k,j}$, so $\{\Np_k\ge j\}\in\mathcal F^{k,j}$;
and $\Xp_{k,j}$ is independent of $\mathcal F^{k,j}$ under $\Pp_{\nu'}$ (product measure). For
nonnegative $\varphi$, by Tonelli,
\begin{align*}
\E_{\nu'}\Big[\sum_{j=1}^{\Np_k}\varphi(\Xp_{k,j})\Big]
&=\sum_{j\ge1}\E_{\nu'}\big[\ind\{\Np_k\ge j\}\varphi(\Xp_{k,j})\big]\\
&=\sum_{j\ge1}\Pp_{\nu'}(\Np_k\ge j)\int\varphi\,d\nu'_k && \text{(independence)}\\
&=\E_{\nu'}[\Np_k]\int\varphi\,d\nu'_k,
\end{align*}
where the product of two independent nonnegative variables has the product of the expectations
(also with the value $\infty$ and $0\cdot\infty=0$). For a real $\varphi$, apply this to
$\varphi^+$ and $\varphi^-$; both terms are finite and the difference is the identity.
\end{proof}

\subsection{Proof of Proposition~\ref{prop:source-kl}}\label{sec:proof-source-kl}

\begin{proof}
\uses{lem:source-law, lem:wald, lem:kl-neg, def:laws, def:source-space, def:size-rule, def:source-data}
\emph{Idea.} By Lemma~\ref{lem:source-law}, the law of $\Dd$ has a density, the product of the
per-sample likelihood ratios. Wald's identity (Lemma~\ref{lem:wald}) computes the expectation of its
logarithm, once the negative part is shown to be integrable.

\emph{Setting.} Write $P:=\PD_{\nu',\Np}$, $\tilde P:=\PD_{\tilde\nu',\Np}$, and let
\begin{equation*}
A:=\{k:\Pp_{\nu'}(\Np_k\ge1)>0\}=\{k:\E_{\nu'}[\Np_k]>0\}.
\end{equation*}
Arms $k\notin A$ contribute $0$ to the right-hand side. \emph{Goal.} Show that
\begin{equation*}
\KL(P,\tilde P)=\sum_{k\in A}\E_{\nu'}[\Np_k]\KL(\nu'_k,\tilde\nu'_k).
\end{equation*}

\emph{Case 1: $\nu'_k\not\ll\tilde\nu'_k$ for some $k\in A$.} Take a Borel set $G$ with
$\tilde\nu'_k(G)=0<\nu'_k(G)$, and let $F(s):=\#\{j\le n_k(s):x_{k,j}\in G\}$, a measurable function
on $\Sd$. By Lemma~\ref{lem:wald} with $\varphi=\ind_G$,
\begin{equation*}
\int F\,dP=\E_{\nu'}[\Np_k]\,\nu'_k(G)>0,\qquad \int F\,d\tilde P=\E_{\tilde\nu'}[\Np_k]\,\tilde\nu'_k(G)=0.
\end{equation*}
So $\tilde P(F\ge1)=0<P(F\ge1)$, $P\not\ll\tilde P$, and both sides equal $\infty$ (the right-hand
side contains $\E_{\nu'}[\Np_k]\cdot\infty$ with $\E_{\nu'}[\Np_k]>0$).

\emph{Case 2: $\nu'_k\ll\tilde\nu'_k$ for every $k\in A$.} Let $r_k$ be a finite version of
$d\nu'_k/d\tilde\nu'_k$ for $k\in A$, and on $\Sd$ let
\begin{equation*}
R(n,x):=\prod_{k\in A}\prod_{j\le n_k}r_k(x_{k,j})\ \text{ if }n_k=0\text{ for every }k\notin A,
\qquad R(n,x):=0\ \text{ otherwise}.
\end{equation*}
\emph{Step 1: $dP=R\,d\tilde P$, slice by slice.} If $n_k\ge1$ for some $k\notin A$, then
$P(\{n\}\times\Sd_n)\le \Pp_{\nu'}(\Np_k\ge1)=0$. Otherwise, $\Bx(n)$ only contains pairs $(k,j)$ with
$k\in A$, where $\bigotimes_{\Bx(n)}\nu'_k$ has density $\prod_{\Bx(n)}r_k(x_{k,j})$ with respect to
$\bigotimes_{\Bx(n)}\tilde\nu'_k$, so by Lemma~\ref{lem:source-law} (the same $g_n$ for both laws),
\begin{equation*}
P(\{n\}\times G)=\int_Gg_n\,R(n,\cdot)\,d\bigotimes_{\Bx(n)}\tilde\nu'_k=\int_{\{n\}\times G}R\,d\tilde P.
\end{equation*}
Summing over the countably many slices, $P\ll\tilde P$ with density $R$. Then $R>0$ $P$-almost
surely, and $\KL(P,\tilde P)=\int\log R\,dP$.

\emph{Step 2: the split.} Write $\ell_k:=\log r_k$ and, $P$-almost surely,
\begin{equation*}
\log R=\sum_{k\in A}\sum_{j\le n_k}\ell_k(x_{k,j})=F^+-F^-,\qquad
F^\pm:=\sum_{k\in A}\sum_{j\le n_k}\ell_k^\pm(x_{k,j})\ge0.
\end{equation*}
Integrating over $P$ means taking $\E_{\nu'}$ of the same functions of $\Dd$, so by
Lemma~\ref{lem:wald},
\begin{equation*}
\int F^\pm\,dP=\sum_{k\in A}\E_{\nu'}[\Np_k]\int\ell_k^\pm\,d\nu'_k.
\end{equation*}
We now split $\int\log R\,dP$ into these two integrals:
\begin{align}
\int F^-\,dP&=\sum_{k\in A}\E_{\nu'}[\Np_k]\int\ell_k^-\,d\nu'_k
&&\text{(Wald, above)}\notag\\
&\le\frac1e\sum_{k\in A}\E_{\nu'}[\Np_k]<\infty
&&\text{(Lemma~\ref{lem:kl-neg}, hypothesis)},\label{eq:Fminus}\\
(\log R)^++F^-&=F^++(\log R)^-\quad P\text{-a.s.}
&&\text{(both are $\log R+F^-+(\log R)^-$)},\label{eq:parts}\\
\int(\log R)^+dP+\int F^-dP&=\int F^+dP+\int(\log R)^-dP
&&\text{(integrate~\eqref{eq:parts})},\notag\\
\int\log R\,dP&=\int F^+\,dP-\int F^-\,dP
&&\text{(subtract finite terms)}.\notag
\end{align}
In the third line all terms are nonnegative. The terms subtracted in the last line are finite:
$\int F^-dP$ by~\eqref{eq:Fminus}, and $\int(\log R)^-dP\le1/e$ by Lemma~\ref{lem:kl-neg}, applied
to $P\ll\tilde P$. By the two
Wald identities above and~\eqref{eq:Fminus}, the last difference is
\begin{align*}
\int\log R\,dP&=\sum_{k\in A}\E_{\nu'}[\Np_k]\Big(\int\ell_k^+d\nu'_k-\int\ell_k^-d\nu'_k\Big)
&&\text{(each $\E_{\nu'}[\Np_k]\int\ell_k^-d\nu'_k$ is finite)}\\
&=\sum_{k\in A}\E_{\nu'}[\Np_k]\KL(\nu'_k,\tilde\nu'_k)
&&\text{(Definition~\ref{def:laws}, $\nu'_k\ll\tilde\nu'_k$)},
\end{align*}
in $[0,\infty]$. Since $\KL(P,\tilde P)=\int\log R\,dP$ by Step~1, this is the identity of the
proposition in Case~2, and Case~1 was settled above: the proposition holds in both cases.
\end{proof}

\subsection{Proof of Lemma~\ref{lem:mixture}}\label{sec:proof-mixture}

\begin{proof}
\uses{def:source-data, def:game}
This is the integral against a composition product $\mu\otimes\kappa$: for nonnegative measurable
$f$,
\begin{equation*}
\int f\,d(\mu\otimes\kappa)=\int\!\!\int f(s,h)\,\kappa(s,dh)\,\mu(ds),
\end{equation*}
with $\mu=\PD_{\nu',\Np}$ and $\kappa(s)=\Q^T_{\pi(s),\nu}$. For $f=n_k(s)$ the inner integral is
$n_k(s)$, whose $\PD_{\nu',\Np}$-integral is $\E_{\nu'}[\Np_k]$ since $n_k(\Dd)=\Np_k$.
\end{proof}

\subsection{Proof of Lemma~\ref{lem:lml-history}}\label{sec:proof-lml-history}

\begin{proof}
\uses{def:laws, def:learner, def:game}
By induction on $T$. For $T=0$ both laws are the Dirac mass on the empty history. Assume that the
law of $H_t$ is $\iota_t\#\Q^t_{\sigma,\nu}$. By the two conditional laws of the statement, the law
of $H_{t+1}=(H_t,A_{t+1},Y_{t+1})$ is the composition product
$(\iota_t\#\Q^t_{\sigma,\nu})\otimes(\bar\sigma_{t+1}\otimes\bar\nu)$. The law $\iota_t\#\Q^t_{\sigma,\nu}$
is carried by $\iota_t(\Hs_t)$, on which $c$ is the identity, so
$\bar\sigma_{t+1}(\iota_th)=\sigma_{t+1}(h)$ for $h\in\Hs_t$; and $\bar\nu_a$ is the image of $\nu_a$.
Hence
\begin{equation*}
(\iota_t\#\Q^t_{\sigma,\nu})\otimes(\bar\sigma_{t+1}\otimes\bar\nu)
=\iota_{t+1}\#\big(\Q^t_{\sigma,\nu}\otimes(\sigma_{t+1}\otimes\nu)\big)=\iota_{t+1}\#\Q^{t+1}_{\sigma,\nu},
\end{equation*}
by Definition~\ref{def:game}. This is the claim at $t+1$.

For the reading in LML: the observations take a single value, so conditioning on the history of
LML and the current observation is conditioning on $H_{t-1}$, read through $j_{t-1}$. The policy
of index $t-1$ is then the conditional law $\bar\sigma_t(H_{t-1})$ of $A_t$, and the reward
kernel of the bandit environment $\bar\nu$ is the conditional law $\bar\nu_{A_t}$ of $Y_t$. So the
two conditional laws of the statement hold, the first part applies, and the history of $T$ rounds
of LML is $j_T^{-1}(H_T)$, with law $j_T^{-1}\#(\iota_T\#\Q^T_{\sigma,\nu})$.
\end{proof}

\subsection{Proof of Lemma~\ref{lem:online-kl}}\label{sec:proof-online-kl}

\begin{proof}
\uses{def:laws, def:learner, def:game, lem:mixture, lem:lml-history}
Fix $s$ and let $\sigma:=\pi(s)$. Let $\bar\sigma$, $\bar\nu$, $\bar{\tilde\nu}$ and $\iota_T$ be as in
Lemma~\ref{lem:lml-history}, with the bijections $j_t$ of its second part, and take
algorithm--environment sequences of LML for the algorithm built there from $\bar\sigma$, against
the bandit environments $\bar\nu$ and $\bar{\tilde\nu}$ (they exist: LML constructs them as
trajectory measures\leanfn{Learning.trajMeasure}). The divergence decomposition of bandit histories
\citep[(8)]{GMS19}, in its formal version in LML\leanfn{Learning.IsAlgEnvSeq.klDiv_map_history_bandit},
for one algorithm against two stationary environments with rewards in $\mathbb R$, gives for the
laws of the histories of the first $T$ rounds
\begin{equation*}
\KL\big(j_T^{-1}\#(\iota_T\#\Q^T_{\sigma,\nu}),\,j_T^{-1}\#(\iota_T\#\Q^T_{\sigma,\tilde\nu})\big)
=\sum_{t\le T}\E_{\sigma,\nu}\big[\KL(\bar\nu_{A_t},\bar{\tilde\nu}_{A_t})\big],
\end{equation*}
where the laws are identified by Lemma~\ref{lem:lml-history} (\emph{the online law as an LML
history}). The maps $j_T^{-1}$ and $\iota_T$ are injective with measurable left inverses ($j_T$ and
the clamp $c$), so they leave $\KL$ unchanged: $\KL(f\#P,f\#Q)=\KL(P,Q)$ for such a map $f$, by data
processing in both directions. Likewise $\KL(\bar\nu_k,\bar{\tilde\nu}_k)=\KL(\nu_k,\tilde\nu_k)$.
Since
$\sum_{t\le T}\KL(\nu_{A_t},\tilde\nu_{A_t})=\sum_kN_k(T)\KL(\nu_k,\tilde\nu_k)$,
\begin{equation*}
\KL\big(\kap^T_{\pi,\nu}(s),\kap^T_{\pi,\tilde\nu}(s)\big)=\sum_k\E_{\pi(s),\nu}[\Nc_k(T)]\KL(\nu_k,\tilde\nu_k)
\qquad\text{in }[0,\infty].
\end{equation*}
Integrate in $s$ against $\PD_{\nu',\Np}$: the constants $\KL(\nu_k,\tilde\nu_k)\in[0,\infty]$ come out
of the integral of the nonnegative functions $s\mapsto\E_{\pi(s),\nu}[\Nc_k(T)]$ (also when a constant
is $\infty$, with $0\cdot\infty=0$), and Lemma~\ref{lem:mixture} (\emph{mixture over the source})
identifies the integrals with $\E_{\PP}[\Nc_k(T)]$.
\end{proof}

\section{Proofs of the basic properties}\label{sec:proofs-properties}

This appendix proves the results of Section~\ref{sec:properties}, in their order.

\subsection{Proof of Lemma~\ref{lem:prior-game}}

\begin{proof}
\uses{def:size-rule, def:prior-game, def:source-space}
\emph{(a)} $\{\Np=n\}$ is $\Om$ or $\emptyset$, both in $\Gc_n$.

\emph{(b)} The maps $A'_r$, $C(r)$, $Y'_r$ are measurable by induction on $r$ (compositions of
measurable maps, and $Y'_r=\sum_{k,j}\ind\{A'_r=k,\,C_k(r)=j\}\Xp_{k,j}$). Let $\Psi(u,x)\in\N^K$ be
the counts of the game run on $\omega=(u,x)\in\Om$, and for $n\in\N^K$ let $x^{(n)}$ be the stack
family equal to $x$ on $\Bx(n)$ and to $0$ outside.

\emph{Claim: if $x$ and $z$ agree on $\Bx(n)$ and $\Psi(u,x)=n$, then $\Psi(u,z)=n$.} By induction on
$r$, the two games have the same first $r$ actions and observations: if this holds up to $r-1$,
round $r$ has the same action $a$ in both games (same $u$, same past), the count $C_a(r)$ is the
same in both games (same past actions), and the observations are $x_{a,C_a(r)}$ and $z_{a,C_a(r)}$;
in the first game $C_a(r)\le C_a(T')=n_a$ and $C_a(r)\ge1$, so $(a,C_a(r))\in\Bx(n)$ and
$x_{a,C_a(r)}=z_{a,C_a(r)}$. Hence the two games coincide and their final counts agree.

\emph{Conclusion.} Applying the claim to $(x,x^{(n)})$ and to $(x^{(n)},x)$ (they agree on $\Bx(n)$),
$\Psi(u,x)=n$ if and only if $\Psi(u,x^{(n)})=n$. So $\{\Np=n\}=\{\omega:\Psi(\Up,\Xp^{(n)})=n\}$, and
$(\Up,\Xp^{(n)})$ is a measurable function of $\Up$ and of the coordinates in $\Bx(n)$; hence
$\{\Np=n\}\in\Gc_n$. Finally $\sum_kN'_k=T'$, so $\Np_k\le T'$.
\end{proof}

\subsection{Proof of Lemma~\ref{lem:lookahead}}

\begin{proof}
\uses{def:size-rule, def:source-space}
$\Nla$ is measurable (a function of $\Xp_{1,2}$) with values in $\{(1,0,\dots,0),(2,0,\dots,0)\}$.
Let $n:=(1,0,\dots,0)$, so $\Gc_n=\sigma(\Up,\Xp_{1,1})$ and $\{\Nla=n\}=\{\Xp_{1,2}>1/2\}$. Every event
of $\Gc_n$ has the form $\{(\Up,\Xp_{1,1})\in G\}$, so contains, with any $\omega$, every $\tilde\omega$
having the same coordinates $\Up$ and $\Xp_{1,1}$. Take $\omega\in\Om$ with $\Xp_{1,2}(\omega)=1$ and
$\tilde\omega$ equal to $\omega$ except $\Xp_{1,2}(\tilde\omega)=0$ (possible since $\Om$ is a
product): $\omega\in\{\Nla=n\}$ and $\tilde\omega\notin\{\Nla=n\}$. So $\{\Nla=n\}\notin\Gc_n$.
\end{proof}

\subsection{Proof of Proposition~\ref{prop:stopped-games}}\label{sec:proof-stopped-games}

\begin{proof}
\uses{def:source-space, def:size-rule, def:prior-game, def:stopped-game, lem:prior-game}
Throughout, the $j$-th observation of arm $k$ in a prior game is $\Xp_{k,j}$
(Definition~\ref{def:prior-game}: $Y'_r=\Xp_{A'_r,C_{A'_r}(r)}$), and $|C(r)|=r$. For
$m\in\N^K$ let $Y_m:=(\Up,(\Xp_{k,j})_{(k,j)\in\Bx(m)})$, with values in the Polish space
$W_m:=[0,1]^{\N}\times[0,1]^{\Bx(m)}$; since $\Gc_m=\sigma(Y_m)$, for a size rule $\Np$ there are
Borel sets $E_m\subset W_m$ with $\{\Np=m\}=\{Y_m\in E_m\}$.

\emph{(a)} The stopped counts $\Np=\sum_r\ind\{\tau=r\}C(r)$ are measurable, with $|\Np|=\tau\le T'$.
Let $\Psi(u,x)\in\N^K$ be the stopped counts at the point $\omega=(u,x)\in\Om$, and $x^{(n)}$ as
in the proof of Lemma~\ref{lem:prior-game}. \emph{Claim: if $x$ and $z$ agree on $\Bx(n)$ and
$\Psi(u,x)=n$, then $\Psi(u,z)=n$.} Let $r_0:=\tau(u,x)$, so $C(r_0)=n$ in the game run on $(u,x)$.
For $r\le r_0$, the observation of round $r$ in that game is $x_{a,C_a(r)}$ with $a=A'_r$ and
$1\le C_a(r)\le C_a(r_0)=n_a$, a coordinate in $\Bx(n)$. By the induction of the proof of
Lemma~\ref{lem:prior-game}, the games run on $(u,x)$ and on $(u,z)$ have the same actions and
observations in rounds $1,\dots,r_0$. Each event $\{\tau=r\}$, $r\le r_0$, is a function of $u$
and of the first $r$ actions and observations ($\{\tau=r\}\in\mathcal F'_r$), so $\tau(u,z)=r_0$,
and the counts at round $r_0$ agree: $\Psi(u,z)=n$. As for Lemma~\ref{lem:prior-game}, applying
the claim to $(x,x^{(n)})$ and to $(x^{(n)},x)$ gives $\{\Np=n\}=\{\Psi(\Up,\Xp^{(n)})=n\}\in\Gc_n$.

\emph{(b), idea.} A game that follows $\Np$ must, at each vector of counts $m$ it reaches, pull an
arm whose size is not yet reached: arm $2$ on the event $A_m$ below, where the size of arm $1$ is
reached and that of arm $2$ is not, and arm $1$ on the symmetric event $B_m$. These two events are
separated by a Borel set of the data seen at $m$ (Step~1). The game pulls
arm $2$ on that set and arm $1$ off it, and stops at the first $m$ at which the event $\{\Np=m\}$
holds (Step~2); Step~3 checks that it stops at $\Np$.

\emph{(b)} Let $K=2$, $\Np$ a size rule with $T':=\sup|\Np|<\infty$, and fix $m\in\N^2$. Let
\begin{equation*}
A_m:=\{\Np_1=m_1,\ \Np_2>m_2\},\qquad B_m:=\{\Np_1>m_1,\ \Np_2=m_2\},
\end{equation*}
the events where the first coordinate, or the second, has reached its final value at $m$ and the
other has not. \emph{Step 1: a Borel set $S_m\subset W_m$ with $A_m\subset\{Y_m\in S_m\}$ and
$B_m\cap\{Y_m\in S_m\}=\emptyset$.} Write a point of $\Om$ as $\omega=(y,a,b)$ with $y=Y_m(\omega)$,
$a=(\Xp_{1,j})_{j>m_1}$ and $b=(\Xp_{2,j})_{j>m_2}$; every triple $(y,a,b)$ is a point of $\Om$.
The event $A_m=\bigcup_{j>m_2}\{\Np=(m_1,j)\}$ is a countable union of events of the
$\Gc_{(m_1,j)}$, which depend on $\omega$ only through $y$ and $b$; so $A_m=\{(y,b)\in A'\}$ for a
Borel set $A'$. Likewise $B_m=\{(y,a)\in B'\}$ for a Borel set $B'$. If some $y$ had a $b$ with
$(y,b)\in A'$ and an $a$ with $(y,a)\in B'$, the point $(y,a,b)$ would be in $A_m\cap B_m=\emptyset$.
So the projections of $A'$ and $B'$ on $W_m$ are disjoint. They are analytic sets (projections of
Borel subsets of Polish spaces), so by the Lusin separation theorem \citep[Theorem~14.7]{K95} there
is a Borel set $S_m\subset W_m$ that contains the first and is disjoint from the second; it has the
two properties.

\emph{Step 2: the game.} For a history $h=(a_i,y_i)_{i<r}$ and $u\in[0,1]^{\N}$, let $c(h)\in\N^2$
be the counts of $h$ and $w(u,h)\in W_{c(h)}$ the point $(u,(x_{k,j}))$ where $x_{k,j}$ is the
$j$-th reward of arm $k$ in $h$; both are measurable. Let
\begin{equation*}
\alpha_r(u,h):=2\ \text{if }w(u,h)\in S_{c(h)},\qquad\alpha_r(u,h):=1\ \text{otherwise},
\end{equation*}
for $r=1,\dots,T'$, and $\tau:=\min\{r\le T':\,w(\Up,H'_r)\in E_{C(r)}\}$ (with $\min\emptyset:=T'$),
where $H'_r:=(A'_i,Y'_i)_{i\le r}$. Then $\{\tau=r\}\in\mathcal F'_r$, and in the game
$w(\Up,H'_r)=Y_{C(r)}$, because the $j$-th observation of arm $k$ is $\Xp_{k,j}$.

\emph{Step 3: the game follows $\Np$.} Fix $\omega$. We show by induction that $C(r)\le \Np(\omega)$
for $r\le|\Np(\omega)|$. It holds at $r=0$. Let $r<|\Np(\omega)|$ and $m:=C(r)\le \Np(\omega)$, so
$m\ne \Np(\omega)$ since $|m|=r$. If $\omega\in A_m$, then $Y_m(\omega)\in S_m$, the game pulls arm $2$
and $m+e_2\le \Np(\omega)$ because $\Np_2(\omega)>m_2$; if $\omega\in B_m$, then $Y_m(\omega)\notin S_m$,
the game pulls arm $1$ and $m+e_1\le \Np(\omega)$; otherwise $\Np(\omega)>m$ in both coordinates and
either pull keeps $C(r+1)\le \Np(\omega)$. So for $r<|\Np(\omega)|$, $C(r)\ne \Np(\omega)$, that is
$Y_{C(r)}(\omega)\notin E_{C(r)}$, while $C(r)=\Np(\omega)$ at $r=|\Np(\omega)|\le T'$. Hence
$\tau(\omega)=|\Np(\omega)|$ and the stopped counts at $\omega$ are $C(\tau)=\Np(\omega)$.

\emph{(c)} The three events of the first three cases are disjoint: the first and second disagree on
$b_2$, the second and third on $b_3$, the first and third on $b_1$. Hence $\Ncyc$ is well defined,
measurable and $|\Ncyc|\le3$. It is a size rule: $\{\Ncyc=e_1+e_2\}=\{b_1=1,b_2=0\}\in\Gc_{e_1+e_2}$,
and likewise for $e_2+e_3$ and $e_3+e_1$; $\{\Ncyc=e_1+e_2+e_3\}$ is the complement of the three
events, all in $\Gc_{e_1+e_2+e_3}$; the other events $\{\Ncyc=n\}$ are empty.

Suppose that a prior game over $T'$ rounds and a stopping time $\tau$ have stopped counts $\Np$ with
the law of $\Ncyc$ under $\Pp_{\nu'}$ for every $\nu'\in\Pc^K$. Let $E_0:=\{\tau=0\}$ and, for
$k\in[K]$, $E_k:=\{\tau\ge1,\ A'_1=k\}$. These events partition $\Om$ and belong to
$\mathcal F'_0=\sigma(\Up)$ (since $A'_1=\alpha_1(\Up)$), and $\Up$ has the law $\lambda^{\otimes\N}$
under every $\Pp_{\nu'}$; so $q_k:=\Pp_{\nu'}(E_k)$ does not depend on $\nu'$, and
$\sum_{k=0}^Kq_k=1$. On $E_0$, $\Np=0$; on $E_k$, $\Np_k\ge C_k(1)=1$. Now:
\begin{itemize}
\item $\Ncyc\ne0$ everywhere, so $q_0\le \Pp_{\nu'}(\Np=0)=\Pp_{\nu'}(\Ncyc=0)=0$.
\item Take $\nu'$ with $\nu'_2=\delta_1$ and $\nu'_3=\delta_0$. Then $b_2=1$ and $b_3=0$ almost
  surely, so $\Ncyc=e_2+e_3$ and $\Ncyc_1=0$ almost surely; hence $q_1\le \Pp_{\nu'}(\Np_1\ge1)
  =\Pp_{\nu'}(\Ncyc_1\ge1)=0$.
\item In the same way, $\nu'_3=\delta_1$, $\nu'_1=\delta_0$ gives $\Ncyc=e_3+e_1$ almost surely and
  $q_2=0$; and $\nu'_1=\delta_1$, $\nu'_2=\delta_0$ gives $\Ncyc=e_1+e_2$ almost surely and $q_3=0$.
\item For $k\ge4$, $\Ncyc_k=0$ everywhere, so $q_k\le \Pp_{\nu'}(\Ncyc_k\ge1)=0$.
\end{itemize}
So $\sum_{k=0}^Kq_k=0$, a contradiction.
\end{proof}

\subsection{Proof of Lemma~\ref{lem:learner-exclusion}}

\begin{proof}
\uses{def:source-data, def:learner, def:source-space, def:size-rule}
Here $\Dd=((1,0,\dots,0),(\Xp_{1,1}))$ does not depend on $\Xp_{1,2}$: for $\omega,\tilde\omega$ equal
except $\Xp_{1,2}(\omega)=1$, $\Xp_{1,2}(\tilde\omega)=0$, $f(\Dd(\omega))=f(\Dd(\tilde\omega))$ while the
targets differ. For the learner: $f(s):=\pi_1(s,\{1\})$ would be measurable with
$f(\Dd)=\ind\{\Xp_{1,2}>1/2\}$ on $\Om$, impossible by the same two points. (The statement is
pointwise on $\Om$; under $\Pp_{\nu'}$ with $\nu'_1$ a Dirac mass, $\Xp_{1,2}$ is almost surely
constant and the identity can hold almost surely.)
\end{proof}

\subsection{Proof of Proposition~\ref{prop:drift}}

\begin{proof}
\uses{def:drift, def:instance, def:laws}
Means of laws on $[0,1]$ lie in $[0,1]$, so $|\mean_k-\mean'_k|\le1$. We check the items in order.
\begin{itemize}
\item (a) $\mean_k-\mean'_k\le|\mean_k-\mean'_k|<L_k$, and $|\mean_k-\mean'_k|<L_k$ gives $\le$.
\item (b) $\mean_k-\mean'_k\le|\mean_k-\mean'_k|\le L_k<L'_k$.
\item (c), first claim. If $L\ge\one$, $\PcL$ is everything (by $|\mean_k-\mean'_k|\le1$), so it
  contains $\{\Drift_L\}$. If $L_i<1$ for some $i$, the pair with $\nu_i=\delta_0$, $\nu'_i=\delta_1$
  and $\nu_k=\nu'_k$ for $k\ne i$ satisfies $\Drift_L$ ($-1<L_i$) and is not in $\PcL$ ($1>L_i$).
\item (c), second claim. If $L_k>1$ for all $k$, $\Drift_L$ holds for every pair
  ($\mean_k-\mean'_k\le1<L_k$). If $L_i\le1$ for some $i$, the pair with $\nu_i=\delta_{L_i}$,
  $\nu'_i=\delta_0$ and $\nu_k=\nu'_k$ otherwise is in $\PcL$ (with $L_i\in(0,1]$, $\delta_{L_i}\in\Pc$)
  and $\mean_i-\mean'_i=L_i\not<L_i$.
\item (d) The ``if'' parts were shown in (c); for the ``only if'' parts, the two pairs of (c) are
  counterexamples (the first is outside $\PcL$ when $L_i<1$, the second violates $\Drift_L$ when
  $L_i\le1$).
\item (e) $\mean_k-\mean'_k=0<L_k$; and $1-0=1\le1$, $1\not<1$. \qedhere
\end{itemize}
\end{proof}

\subsection{Proof of Proposition~\ref{prop:classes}}

\begin{proof}
\uses{def:consistency, def:drift, prop:drift}
\emph{(a)} $\PcL\subset\PcL[L']$, and the other quantifiers of Definition~\ref{def:consistency} do
not depend on $L$, so a condition required on $\PcL[L']$ is required on the subset $\PcL$.

\emph{(b)} For $L\ge\one$, $\PcL=\Pc^K\times\Pc^K=\PcL[\one]$ by Proposition~\ref{prop:drift}(d), so
$\ClL=\ClL[\one]$. For every $L$, $L\vee\one\ge L$, so by (a) and the first claim,
$\ClL[\one]=\ClL[L\vee\one]\subset\ClL$; hence $\ClL[\one]\subset\bigcap_L\ClL\subset\ClL[\one]$.
\end{proof}

\subsection{Proof of Proposition~\ref{prop:witness}}

\begin{proof}
\uses{def:consistency, def:learner, def:game, def:instance, lem:mixture, prop:classes, lem:lml-history}
\emph{The source-ignoring learner.} For every size map, $\pis(s)=\sigma$ for all $s$, so
Lemma~\ref{lem:mixture} gives
\begin{equation*}
\E_{\nu',\rhos,\nu,\pis}[\Nc_k(T)]=\int\E_{\sigma,\nu}[\Nc_k(T)]\,\PD(ds)=\E_{\sigma,\nu}[\Nc_k(T)],
\end{equation*}
which is $o(T^\alpha)$ for every suboptimal $k$ and $\alpha>0$, for every pair, whatever $\nu'$ and
$\rhos$: $\pis\in\ClL[\one]$, and $\ClL[\one]\subset\ClL$ by Proposition~\ref{prop:classes}(b).

\emph{UCB.} Let $\sigma$ be UCB as in the statement, on histories with rewards in $[0,1]$, and let
$\nu\in\Pc^K$. The UCB algorithm of LML acts on histories with real rewards; on histories with
rewards in $[0,1]$ it chooses its arm by the rule of the statement, so it agrees there with
$\bar\sigma=\sigma\circ c$ of Lemma~\ref{lem:lml-history}, since the clamp $c$ does not change such
histories. Against the bandit environment $\bar\nu$ the rewards are in $[0,1]$ almost surely, so
the conditional laws of the actions of a run of LML's UCB against $\bar\nu$ agree almost surely with
those given by $\bar\sigma$: the run satisfies the two conditional laws of
Lemma~\ref{lem:lml-history} (\emph{the online law as an LML history}), and the law of its first $n$
rounds is $\iota_n\#\Q^n_{\sigma,\nu}$. Hence the expected numbers of pulls of the run are
$\E_{\sigma,\nu}[\Nc_a(n)]$. A law on $[0,1]$ is sub-Gaussian with variance proxy $1/4$ around its
mean (Hoeffding's lemma, in Mathlib\leanfn{ProbabilityTheory.hasSubgaussianMGF_of_mem_Icc}), and
for $c>2\cdot\frac14$ the regret bound of UCB, due to \citet{ACF02} and formalised in LML with the
following constants\leanfn{Bandits.UCB.regret_le_of_gt_two} (variance proxy $1/4$), gives, for
every $n\in\N$,
\begin{equation*}
\sum_a\gap_a\E_{\sigma,\nu}[\Nc_a(n)]\le\sum_{a:\,\gap_a>0}\Big(\frac{8c\log(n+1)}{\gap_a}
+\gap_a\Big(4+\frac{1/2}{c-1/2}\Big)\Big).
\end{equation*}
Each term on the left is nonnegative, so $\E_{\sigma,\nu}[\Nc_k(n)]=O(\log n)$ for every suboptimal
$k$, which is $o(n^\alpha)$ for every $\alpha>0$.
\end{proof}

\subsection{Proof of Proposition~\ref{prop:exclusion}}

\begin{proof}
\uses{def:consistency, def:learner, def:source-data, def:drift, def:game, def:instance, def:laws, lem:mixture, lem:prior-game}
$\pitr$ is a transfer learner: its kernels are Dirac masses at a measurable function of $s$ ($\hat k$
is measurable on each slice $\Sd_n$). Take $\nu_1=\nu'_1=\Ber(2/3)$, $\nu_2=\nu'_2=\Ber(1/3)$ and
$\nu_k=\nu'_k=\delta_0$ for $k\ge3$: then $(\nu,\nu')\in\PcL$ for every $L$ (no drift), and arm $2$ is
suboptimal. Take the source strategy $\Np^{(T)}\equiv\one$ (a size rule by
Lemma~\ref{lem:prior-game}(a)). On the event $\{\Xp_{1,1}=0,\,\Xp_{2,1}=1\}$, of probability
$\frac13\cdot\frac13=\frac19$, the source empirical means are $0,1,0,\dots,0$, so $\hat k(\Dd)=2$ and
$\Nc_2(T)=T$. By Lemma~\ref{lem:mixture},
\begin{equation*}
\E_{\nu',\rhos,\nu,\pitr}[\Nc_2(T)]\ge T\,\PD(\hat k=2)\ge\frac T9,\qquad\text{so}\qquad
\frac{\E_{\nu',\rhos,\nu,\pitr}[\Nc_2(T)]}{T^{1/2}}\ge\frac{T^{1/2}}9\not\to0. \qedhere
\end{equation*}
\end{proof}

\section{Sources of the techniques}\label{sec:attribution}

For each technique, where it is used and the earliest source we found.

\begin{center}
\begin{tabular}{@{}p{0.45\textwidth}p{0.49\textwidth}@{}}
\toprule
Technique (where) & Earliest source found \\
\midrule
Source model: prior counts chosen by an unknown, possibly adaptive strategy, independent of the
target (Definitions~\ref{def:size-rule}--\ref{def:game}) & \citet[Section~3.2]{PMRM26} \\
The class $\ClL$ (Definition~\ref{def:consistency}) & \citet[Section~4.2]{PMRM26}, copied \\
Size rules as stopping points of a multiparameter filtration (Definition~\ref{def:size-rule}) &
bandits as control over a multiparameter filtration: \citet{M86}, \citet{KM98}; see also
\citet{CW75} \\
Counts of a bandit game, stopped or not, are a stopping point (Lemma~\ref{lem:prior-game},
Proposition~\ref{prop:stopped-games}(a)) & \citet[Section~1.3]{MV81} (the end point of a policy is
a predictable stopping point); for bandits, \citet{KM98}; see also \citet{KS81}; the induction is
written here \\
Every stopping point of two parameters is reached by a policy (Proposition~\ref{prop:stopped-games}(b))
& \citet[Theorem~4]{MV81}, for one probability measure under $(\mathrm{F}4)$, by separating the
events $A_m$ and $B_m$; here the separation is pointwise, by the Lusin separation theorem
\citep[Theorem~14.7]{K95} \\
A stopping point of three parameters that no policy reaches (Proposition~\ref{prop:stopped-games}(c))
& the cyclic pattern and the argument (no argument is read first on every input): Berry's stable
non-sequential function \citep[Proposition~II.2.2]{Ber76}, see \citet{Cur03}; for three independent
Markov chains, \citet[Section~5.2]{MV81}. In law and with randomisation, for mechanisms missing at
random: \citet{RG97}, see also \citet{GR97}. Written here for size rules \\
Factorisation of the law of the source data (Lemma~\ref{lem:source-law}) & ignorability under
missing at random: \citet{R76}; \citet[Theorem~1]{SGJC13}; \citet[Theorem~1]{FDS18}; see also
\citet{BW88} \\
Wald's identity at a stopping point (Lemma~\ref{lem:wald}) & \citet{W44}; for the log-likelihood
ratio of a bandit at a stopping time, \citet[Appendix~A, (18)]{KCG16} \\
Source term $\E[\Np_k]\KL(\nu'_k,\tilde\nu'_k)$ (Proposition~\ref{prop:source-kl}) & for bandit-game
sources, \citet[Lemma~22]{PMRM26}; offline data on the same instance,
\citet[Appendix~B.1, (10)]{AJSS23} \\
Divergence decomposition of bandit histories (Lemma~\ref{lem:online-kl}) & \citet[Lemma~1]{KCG16};
as an identity, \citet[(8)]{GMS19}; formalised in LML \\
Chain rule for composition products (Theorem~\ref{thm:main}) & textbook; formalised in LML \\
Data processing (Corollary~\ref{cor:fundamental}) & \citet[Lemma~1]{GMS19} \\
Look-ahead counterexample (Theorem~\ref{thm:main}(b)) & an instance of a mechanism that is not
missing at random \citep{R76} \\
The learner that follows the source (Proposition~\ref{prop:exclusion}) & source not known \\
Regret bound of UCB (Proposition~\ref{prop:witness}) & \citet{ACF02}; formalised in LML, with the
constants used here \\
\bottomrule
\end{tabular}
\end{center}

\section{Notation}\label{sec:notation}

Every symbol is a link to its definition.

\begin{center}
\begin{tabular}{@{}p{0.27\textwidth}p{0.47\textwidth}p{0.2\textwidth}@{}}
\toprule
Symbol & Meaning & Defined in \\
\midrule
$\Pc$, $\Ber(p)$ & laws on $[0,1]$, Bernoulli law & Definition~\ref{def:laws} \\
$\KL$, $\kl$ & Kullback--Leibler divergence, Bernoulli divergence & Definition~\ref{def:laws} \\
$\mean_k$, $\mean'_k$, $\mean^\star$, $\gap_k$ & target and source means, best mean, gaps & Definition~\ref{def:instance} \\
$\Om$, $\Up$, $\Xp_{k,j}$, $\Pp_{\nu'}$ & source space, randomisation, stacks, their law & Definition~\ref{def:source-space} \\
$\Bx(n)$, $\Gc_n$ & box of $n$, its $\sigma$-algebra & Definition~\ref{def:source-space} \\
$\Np$, $\rhos$ & size map or size rule, source strategy & Definition~\ref{def:size-rule} \\
$\Sd$, $\Sd_n$, $\Dd$, $\PD_{\nu',N'}$ & space of source data, its slices, source data, their law & Definition~\ref{def:source-data} \\
$\mhat_k(s)$ & source empirical mean & Definition~\ref{def:source-data} \\
$\sy{s:game0}{A'_r}$, $\sy{s:game0}{Y'_r}$, $\sy{s:game0}{C_k(r)}$ & actions, observations and counts of a prior game & Definition~\ref{def:prior-game} \\
$\sy{s:stopped}{\mathcal F'_r}$, $\sy{s:stopped}{\tau}$ & filtration of a prior game, stopping time & Definition~\ref{def:stopped-game} \\
$\Hs_t$, $\sy{s:Hs}{\sigma}$, $\sy{s:Hs}{\pi}$, $\pis$ & histories, bandit algorithm, transfer learner, source-ignoring learner & Definition~\ref{def:learner} \\
$\Q^T_{\sigma,\nu}$, $\kap^T_{\pi,\nu}$, $\PP^T_{\nu',N',\nu,\pi}$ & law of the online history, its kernel, law of the game & Definition~\ref{def:game} \\
$\Nc_k(T)$ & number of pulls of arm $k$ in the first $T$ rounds & Definition~\ref{def:game} \\
$\Drift_L$, $\PcL$, $\one$ & drift condition, closed two-sided set, vector of ones & Definition~\ref{def:drift} \\
$\ClL$ & consistency class & Definition~\ref{def:consistency} \\
$\Nla$ & look-ahead map & Theorem~\ref{thm:main} \\
$\sy{s:ek}{e_k}$, $\sy{s:ek}{|n|}$, $\Ncyc$ & unit vector, total size, cyclic size rule & Proposition~\ref{prop:stopped-games} \\
$\pitr$ & learner that follows the source & Proposition~\ref{prop:exclusion} \\
\bottomrule
\end{tabular}
\end{center}

\bibliographystyle{plainnat}
\bibliography{refs,../../papers/library}

\end{document}
